\documentclass{article}%
\usepackage[T1]{fontenc}%
\usepackage[utf8]{inputenc}%
\usepackage{lmodern}%
\usepackage{textcomp}%
\usepackage{lastpage}%
\usepackage{geometry}%
\geometry{margin=1in,headheight=14pt,headsep=25pt}%
\usepackage{hyperref}%
\usepackage{enumitem}%
\usepackage{fancyhdr}%
\usepackage{xcolor}%
\usepackage{url}%
\usepackage{breakurl}%
\usepackage{amsmath}%
\usepackage{amssymb}%
\usepackage{mathtools}%
\usepackage{amsthm}%
\usepackage{thmtools}%
\usepackage{float}%
\usepackage{subcaption}%
\usepackage[font=small,labelfont=bf]{caption}%
%

            \hypersetup{
                colorlinks=true,
                linkcolor=blue,
                filecolor=magenta,
                urlcolor=blue,
                breaklinks=true,
                bookmarks=true,
                pdfborder={0 0 0}
            }
            \urlstyle{same}
            \Urlmuskip=0mu plus 1mu  % URL breaking in nicer spots

            % Math configuration
            \everymath{\displaystyle}
            \setlength{\jot}{10pt}  % Increase spacing between equations

            \pagestyle{fancy}
            \fancyhf{}
            \rhead{Generated on \today}
            \cfoot{\thepage}

            % Better Unicode support
            \usepackage[utf8]{inputenc}
            \usepackage{newunicodechar}
            \setlength{\itemsep}{1em} % Adds spacing between items

            % Configure itemize settings globally
            \setlist[itemize]{nosep}
            \setlist[itemize]{leftmargin=*}

            % Remove section numbering
            \renewcommand{\thesection}{}
            \renewcommand{\thesubsection}{}

            % Adjust caption spacing around figures
            \setlength{\abovecaptionskip}{0pt}
            \setlength{\belowcaptionskip}{5pt}
        %
\lhead{Lecture: 2024{-}08{-}27{-}ExSimplex}%
\title{2024{-}08{-}27{-}ExSimplex}%
\author{Generated by Scribe.AI}%
\date{\today}%
%
\begin{document}%
\normalsize%
\maketitle%
\section{Slide 1}%
\label{sec:Slide1}%
\subsection{Content}%

\text{Simplex Method [LV] p.20} \\
\underline{\text{max }} \quad \zeta = 3x_1 + 2x_2 \\
\text{Subject to} \quad -x_1 + 3x_2 \le 12 \\
\quad x_1 + x_2 \le 8 \\
\quad 2x_1 - x_2 \le 10 \\
\quad x_1, x_2 \ge 0 \\
\text{Slack variable form:} \\
w_1 = 12 + x_1 - 3x_2 \\
w_2 = 8 - x_1 - x_2 \\
w_3 = 10 - 2x_1 + x_2 \\
x_1, x_2, x_3, w_1, w_2, w_3 \ge 0
%
\subsection{Description}%
Slide 1 of 44 in a Linear Programming course introduces the simplex method by presenting a linear programming problem. The problem is to maximize the objective function $\zeta = 3x_1 + 2x_2$ subject to three inequality constraints: $-x_1 + 3x_2 \le 12$, $x_1 + x_2 \le 8$, and $2x_1 - x_2 \le 10$.  Non-negativity constraints $x_1, x_2 \ge 0$ are also included. The slide then shows how to convert this problem into its slack variable form by introducing slack variables $w_1, w_2, w_3 \ge 0$.  The resulting equations are: $w_1 = 12 + x_1 - 3x_2$, $w_2 = 8 - x_1 - x_2$, and $w_3 = 10 - 2x_1 + x_2$. This transformation is a standard step in preparing a linear program for solution using the simplex method. The reference "[LV] p.20" likely indicates a textbook or other source where this example can be found. The context of the course is crucial because this slide lays the groundwork for the simplex method, a fundamental algorithm in linear programming.%
\subsection{Figures}%
\begin{figure}[H]%
\centering%
\begin{subfigure}[b]{0.15\textwidth}%
\centering%
\includegraphics[width=\textwidth]{/Users/ashoksaravanan/Coding/ScribeLec/server/output/MA 421/lectures/2024-08-27-ExSimplex/figures/1.1.png}%
\caption{Handwritten notes showing a linear programming problem in standard form and its equivalent slack variable form. The objective function is to maximize $\zeta = 3x_1 + 2x_2$, subject to three inequality constraints and non-negativity constraints on the variables $x_1$ and $x_2$. Slack variables $w_1$, $w_2$, and $w_3$ are introduced to convert the inequalities into equalities. The resulting system of equations is the slack variable form of the linear program.}%
\end{subfigure}%
\end{figure}

%
\newpage%
\section{Slide 2}%
\label{sec:Slide2}%
\subsection{Content}%

\text{2}
%
\subsection{Description}%
Slide 2 of 44 in a Linear Programming course illustrates the graphical representation of a linear inequality. A straight line representing the equation $ax + by = 0$ is plotted on the $x_1$-$x_2$ plane.  Arrows indicate the regions where $ax + by < 0$ and $ax + by > 0$. The axes are labeled $x_1 = 0$ and $x_2 = 0$, indicating that the graph shows the feasible region in the non-negative orthant. This slide provides a visual aid to understand how linear inequalities define regions in the solution space of a linear program. The example is simple, but it sets the stage for more complex problems with multiple constraints that will be encountered later in the course.%
\subsection{Figures}%
\begin{figure}[H]%
\centering%
\begin{subfigure}[b]{0.15\textwidth}%
\centering%
\includegraphics[width=\textwidth]{/Users/ashoksaravanan/Coding/ScribeLec/server/output/MA 421/lectures/2024-08-27-ExSimplex/figures/2.1.png}%
\caption{A graph showing a line $ax + by = 0$ in the $x_1$-$x_2$ plane. The regions $ax + by < 0$ and $ax + by > 0$ are indicated with arrows. The axes are labeled $x_1 = 0$ and $x_2 = 0$.}%
\end{subfigure}%
\end{figure}

%
\newpage%
\section{Slide 3}%
\label{sec:Slide3}%
\subsection{Content}%

\text{eg. } ax + by = w, \quad w < 0 \\
ax + by = w, \quad w = 0 \\
ax + by = w, \quad w > 0
%
\subsection{Description}%
Slide 3 of 44 in a Linear Programming course illustrates the graphical representation of a linear inequality and its relationship to the feasible region.  A straight line representing the equation $ax + by = w$ is plotted on the $x_1$-$x_2$ plane. Arrows indicate the regions where $ax + by < w$, $ax + by = w$, and $ax + by > w$. The axes are labeled $x_1 = 0$ and $x_2 = 0$, indicating that the graph shows the feasible region in the non-negative orthant.  The example builds upon the previous slide by showing how different values of $w$ shift the line parallel to itself, affecting the feasible region. This slide helps visualize how linear inequalities define regions in the solution space of a linear program, which is crucial for understanding the simplex method and other graphical solution techniques in linear programming.%
\subsection{Figures}%
\begin{figure}[H]%
\centering%
\begin{subfigure}[b]{0.15\textwidth}%
\centering%
\includegraphics[width=\textwidth]{/Users/ashoksaravanan/Coding/ScribeLec/server/output/MA 421/lectures/2024-08-27-ExSimplex/figures/3.1.png}%
\caption{A graph showing a line $ax + by = w$ in the $x_1$-$x_2$ plane. The regions $ax + by < w$, $ax + by = w$, and $ax + by > w$ are indicated with arrows. The axes are labeled $x_1 = 0$ and $x_2 = 0$.}%
\end{subfigure}%
\end{figure}

%
\newpage%
\section{Slide 4}%
\label{sec:Slide4}%
\subsection{Content}%

%
\subsection{Description}%
Slide 4 of 44 in a Linear Programming course shows a graphical representation of the feasible region for the linear program introduced in slide 1. The feasible region is a pentagon defined by the intersection of five constraints: $x_1 \ge 0$, $x_2 \ge 0$, $-x_1 + 3x_2 \le 12$, $x_1 + x_2 \le 8$, and $2x_1 - x_2 \le 10$. Each line corresponds to one of the constraints, with arrows indicating the direction of the inequality. The feasible region is the area where all constraints are satisfied simultaneously, shaded in orange.  The equations of the lines are written next to the lines themselves, and the corresponding slack variable ($w_i = 0$) is indicated for each constraint. This visual representation is crucial for understanding the concept of a feasible region and how the constraints limit the possible solutions of the linear program.  The slide serves as a visual aid to complement the algebraic representation from the previous slide and sets the stage for the application of the simplex method.%
\subsection{Figures}%
\begin{figure}[H]%
\centering%
\begin{subfigure}[b]{0.15\textwidth}%
\centering%
\includegraphics[width=\textwidth]{/Users/ashoksaravanan/Coding/ScribeLec/server/output/MA 421/lectures/2024-08-27-ExSimplex/figures/4.1.png}%
\caption{A graph showing the feasible region of a linear programming problem. The feasible region is a pentagon bounded by the lines $x_1 = 0$, $x_2 = 0$, $-x_1 + 3x_2 = 12$, $x_1 + x_2 = 8$, and $2x_1 - x_2 = 10$. The feasible region is shaded in orange. Arrows indicate the direction of the inequalities. The axes are labeled $x_1$ and $x_2$.}%
\end{subfigure}%
\end{figure}

%
\newpage%
\section{Slide 5}%
\label{sec:Slide5}%
\subsection{Content}%

%
\subsection{Description}%
Slide 5 of 44 in a Linear Programming course presents a graphical solution to a linear programming problem.  The feasible region is a pentagon, shaded in orange, defined by the intersection of five constraints: $x_1 \ge 0$, $x_2 \ge 0$, $-x_1 + 3x_2 \le 12$, $x_1 + x_2 \le 8$, and $2x_1 - x_2 \le 10$. Each constraint is represented by a line, with arrows indicating the direction of the inequality. The vertices of the feasible region are labeled A, B, C, D, and O, with their coordinates given in terms of the original variables ($x_1$, $x_2$) and slack variables ($w_1$, $w_2$, $w_3$).  The point E lies outside the feasible region and is explicitly labeled as "not feasible".  The slide visually demonstrates the concept of a feasible region and how the optimal solution will lie at one of the vertices of this region. This is a crucial step in understanding the graphical method for solving linear programming problems, which will be further developed in subsequent slides. The slide builds upon the previous slide by showing a concrete example of a feasible region and its vertices.%
\subsection{Figures}%
\begin{figure}[H]%
\centering%
\begin{subfigure}[b]{0.15\textwidth}%
\centering%

\includegraphics[width=\textwidth]{/Users/ashoksaravanan/Coding/ScribeLec/server/output/MA 421/lectures/2024-08-27-ExSimplex/figures/5.1.png}%
\caption{A graph showing the feasible region of a linear programming problem. The feasible region is a pentagon bounded by the lines $x_1 = 0$, $x_2 = 0$, $-x_1 + 3x_2 = 12$, $x_1 + x_2 = 8$, and $2x_1 - x_2 = 10$. The feasible region is shaded in orange. The vertices of the feasible region are labeled A, B, C, D, and O.  The coordinates of the vertices are given as follows: O($x_1$=0, $x_2$=0), A($x_2$=0, $w_3$=0), B($w_2$=0, $w_3$=0), C($w_1$=0, $w_2$=0), D($x_1$=0, $w_1$=0). Point E is shown outside the feasible region and labeled as "not feasible". Arrows indicate the direction of the inequalities. The axes are labeled $x_1$ and $x_2$.  The equations of the lines are written next to the lines themselves, and the corresponding slack variable ($w_i = 0$) is indicated for each constraint.}%
\end{subfigure}%
\end{figure}

%
\newpage%
\section{Slide 6}%
\label{sec:Slide6}%
\subsection{Content}%

(0) \quad \text{max } \zeta = 3x_1 + 2x_2 \\
\text{subject to } \begin{cases}
w_1 = 12 + x_1 - 3x_2 \\
w_2 = 8 - x_1 - x_2 \\
w_3 = 10 - 2x_1 + x_2
\end{cases} \\
x_1, x_2, x_3, w_1, w_2, w_3 \ge 0 \\
x_1, x_2, x_3 \text{ on the RHS} - \text{non-basic variables} \\
w_1, w_2, w_3 \text{ on the LHS} - \text{basic variables} \\
\text{Idea of Simplex Method:} \\
(1) \text{ set all non-basic var. to zero - correspond to a vertex} \\
(2) \text{ choose a new set of (non)basic vars, to improve } \zeta \\
(1) \text{ or } (2) \text{ is the same as moving from vertex to vertex}
%
\subsection{Description}%
Slide 6 of 44 in a Linear Programming course introduces the simplex method.  The slide presents a linear programming problem: maximize $\zeta = 3x_1 + 2x_2$ subject to the constraints $w_1 = 12 + x_1 - 3x_2$, $w_2 = 8 - x_1 - x_2$, and $w_3 = 10 - 2x_1 + x_2$, with all variables non-negative.  The constraints are rewritten using slack variables ($w_1, w_2, w_3$).  The variables $x_1, x_2, x_3$ are identified as non-basic variables (on the right-hand side, RHS), and $w_1, w_2, w_3$ as basic variables (on the left-hand side, LHS). The core idea of the simplex method is explained: (1) setting all non-basic variables to zero corresponds to finding a vertex of the feasible region; (2) iteratively choosing a new set of basic and non-basic variables to improve the objective function $\zeta$.  The slide emphasizes that this iterative process is equivalent to moving from one vertex to another within the feasible region, searching for the optimal solution. This slide bridges the graphical representation of the feasible region (previous slides) with the algebraic approach of the simplex method.%
\subsection{Figures}%
\begin{figure}[H]%
\centering%
\begin{subfigure}[b]{0.15\textwidth}%
\centering%
\includegraphics[width=\textwidth]{/Users/ashoksaravanan/Coding/ScribeLec/server/output/MA 421/lectures/2024-08-27-ExSimplex/figures/6.1.png}%
\caption{Handwritten notes describing a linear programming problem and the idea behind the simplex method. The problem is to maximize $\zeta = 3x_1 + 2x_2$ subject to three constraints. The constraints are written in a form where slack variables $w_1, w_2, w_3$ are explicitly defined. The variables $x_1, x_2, x_3$ are classified as non-basic variables, while $w_1, w_2, w_3$ are basic variables. The idea of the simplex method is explained in two steps: (1) setting non-basic variables to zero to find a vertex of the feasible region, and (2) choosing a new set of basic and non-basic variables to improve the objective function.  The equivalence between these steps and moving from vertex to vertex in the feasible region is highlighted.}%
\end{subfigure}%
\end{figure}

%
\newpage%
\section{Slide 7}%
\label{sec:Slide7}%
\subsection{Content}%

\circled{1} \quad \underline{\text{Initial Step/Guess}} \\
\text{max } \zeta = 3x_1 + 2x_2 \\
\text{subject to } \begin{cases}
w_1 = 12 + x_1 - 3x_2 \\
w_2 = 8 - x_1 - x_2 \\
w_3 = 10 - 2x_1 + x_2
\end{cases} \\
x_1, x_2, x_3, w_1, w_2, w_3 \ge 0 \\
\underline{\text{Set } x_1, x_2 = 0 \text{ (O) } \implies \zeta = 0}
%
\subsection{Description}%
Slide 7 of 44 in a Linear Programming course details the initial step of the simplex method applied to the linear program introduced in previous slides. The slide begins by stating "Initial Step/Guess" and then presents the linear program: maximize $\zeta = 3x_1 + 2x_2$ subject to the constraints $w_1 = 12 + x_1 - 3x_2$, $w_2 = 8 - x_1 - x_2$, and $w_3 = 10 - 2x_1 + x_2$, with $x_1, x_2, x_3, w_1, w_2, w_3 \ge 0$.  The initial step involves setting the non-basic variables $x_1$ and $x_2$ to zero. This corresponds to the origin (O) in the feasible region graph from previous slides.  The resulting value of the objective function $\zeta$ is calculated as 0, which is underlined for emphasis.  The slide sets the stage for the iterative process of the simplex method, where the algorithm will move from this initial vertex to find the optimal solution.%
\subsection{Figures}%
\begin{figure}[H]%
\centering%
\begin{subfigure}[b]{0.15\textwidth}%
\centering%
\includegraphics[width=\textwidth]{/Users/ashoksaravanan/Coding/ScribeLec/server/output/MA 421/lectures/2024-08-27-ExSimplex/figures/7.1.png}%
\caption{Handwritten notes describing the initial step of the simplex method. The problem is to maximize $\zeta = 3x_1 + 2x_2$ subject to three constraints. The constraints are written in a form where slack variables $w_1, w_2, w_3$ are explicitly defined. The variables $x_1, x_2, x_3$ are non-negative, as are the slack variables. The initial step involves setting $x_1$ and $x_2$ to zero, resulting in $\zeta = 0$. The underlined text highlights the initial guess and the resulting objective function value.}%
\end{subfigure}%
\end{figure}

%
\newpage%
\section{Slide 8}%
\label{sec:Slide8}%
\subsection{Content}%

\circled{1} \quad \text{Improve } (O) \ (x_1 = 0, x_2 = 0) \\
\text{increase } x_1, \text{ keep } x_2 = 0 \\
\text{(max coeff. rule)} \\
\text{max } \zeta = 3x_1 + 2x_2 \\
\text{subject to } \begin{cases}
w_1 = 12 + x_1 - 3x_2 \\
w_2 = 8 - x_1 - x_2 \\
w_3 = 10 - 2x_1 + x_2
\end{cases} \\
x_1, x_2, x_3, w_1, w_2, w_3 \ge 0 \\
w_1: x_1 \text{ can be as large as possible} \\
w_2: 8 - x_1 \ge 0 \implies x_1 \le 8 \\
w_3: 10 - 2x_1 \ge 0 \implies x_1 \le 5 \\
\text{Set } x_1 = 5, \implies w_3 = 0 \implies A = (x_2 = 0, w_3 = 0)
%
\subsection{Description}%
Slide 8 of 44 in a Linear Programming course demonstrates the first iteration of the simplex method.  Building upon the previous slide, it aims to improve the initial solution (O) where $x_1 = 0$ and $x_2 = 0$, resulting in $\zeta = 0$. The strategy is to increase $x_1$ while keeping $x_2 = 0$ (the "max coeff. rule," referring to selecting the variable with the largest coefficient in the objective function).  The constraints are examined to find the maximum feasible value of $x_1$.  The inequalities $8 - x_1 \ge 0$ and $10 - 2x_1 \ge 0$ are derived from the constraints $w_2 \ge 0$ and $w_3 \ge 0$ respectively, yielding $x_1 \le 8$ and $x_1 \le 5$. The smaller value, $x_1 = 5$, is chosen to maintain feasibility.  Substituting $x_1 = 5$ and $x_2 = 0$ into the constraints shows that $w_3$ becomes 0, indicating that the new vertex A is reached. The coordinates of A are given as $(x_2 = 0, w_3 = 0)$, representing the new basic feasible solution.  The slide visually connects the algebraic steps of the simplex method with the geometric interpretation of moving from one vertex to another in the feasible region.%
\subsection{Figures}%
\begin{figure}[H]%
\centering%
\begin{subfigure}[b]{0.15\textwidth}%
\centering%
\includegraphics[width=\textwidth]{/Users/ashoksaravanan/Coding/ScribeLec/server/output/MA 421/lectures/2024-08-27-ExSimplex/figures/8.1.png}%
\caption{Handwritten notes describing the first iteration of the simplex method. The problem is to maximize $\zeta = 3x_1 + 2x_2$ subject to three constraints. The constraints are written in a form where slack variables $w_1, w_2, w_3$ are explicitly defined. The variables $x_1, x_2, x_3$ are non-negative, as are the slack variables. The initial step involves setting $x_1$ and $x_2$ to zero, resulting in $\zeta = 0$. The underlined text highlights the initial guess and the resulting objective function value. The notes then describe how to improve the solution by increasing $x_1$ while keeping $x_2 = 0$, based on the maximum coefficient rule. The constraints are analyzed to determine the maximum value of $x_1$ that maintains feasibility, which is found to be 5. Setting $x_1 = 5$ results in $w_3 = 0$, leading to a new vertex A with coordinates $(x_2 = 0, w_3 = 0)$.}%
\end{subfigure}%
\end{figure}

%
\newpage%
\section{Slide 9}%
\label{sec:Slide9}%
\subsection{Content}%

\circled{Ⅱ} \quad A \left( x_2 = 0, w_3 = 0 \right) \\
\underline{\text{Interchange } x_1 \text{ and } w_3 \text{ with }} \\
\underline{w_3 \text{ leaving basic and } x_1 \text{ entering basic}} \\
\zeta = 3x_1 + 2x_2 \\
w_1 = 12 + x_1 - 3x_2 \\
w_2 = 8 - x_1 - x_2 \\
w_3 = 10 - 2x_1 + x_2 \\
\implies \begin{cases}
w_1 = 12 + 5 - \frac{w_3}{2} + \frac{1}{2}x_2 - 3x_2 = 17 - \frac{w_3}{2} - \frac{5}{2}x_2 \\
w_2 = 8 - 5 + \frac{w_3}{2} - \frac{1}{2}x_2 - x_2 = 3 + \frac{w_3}{2} - \frac{3}{2}x_2 \\
\zeta = 3 \left( 5 - \frac{w_3}{2} + \frac{1}{2}x_2 \right) + 2x_2 = 15 - \frac{3}{2}w_3 + \frac{7}{2}x_2
\end{cases}
%
\subsection{Description}%
Slide 9 of 44 in a Linear Programming course shows the second iteration of the simplex method.  It builds on slide 8, where the vertex A ($x_1 = 5, x_2 = 0, w_3 = 0$) was identified.  The slide indicates an interchange of variables: $x_1$ enters the basis, and $w_3$ leaves. This is underlined for emphasis. The original constraints and objective function are shown.  The key step is solving the equation $w_3 = 10 - 2x_1 + x_2$ for $x_1$, resulting in $x_1 = 5 - \frac{w_3}{2} + \frac{1}{2}x_2$. This expression for $x_1$ is then substituted into the other equations to express $w_1$, $w_2$, and $\zeta$ in terms of the new basic variables $w_3$ and $x_2$. The resulting equations are: $w_1 = 17 - \frac{w_3}{2} - \frac{5}{2}x_2$, $w_2 = 3 + \frac{w_3}{2} - \frac{3}{2}x_2$, and $\zeta = 15 - \frac{3}{2}w_3 + \frac{7}{2}x_2$.  This completes the second iteration, preparing for the next step in the simplex algorithm. The slide clearly demonstrates the algebraic manipulation involved in pivoting from one basic feasible solution to another.%
\subsection{Figures}%
\begin{figure}[H]%
\centering%
\begin{subfigure}[b]{0.15\textwidth}%
\centering%
\includegraphics[width=\textwidth]{/Users/ashoksaravanan/Coding/ScribeLec/server/output/MA 421/lectures/2024-08-27-ExSimplex/figures/9.1.png}%
\caption{Handwritten notes describing the second iteration of the simplex method. The problem is to maximize $\zeta = 3x_1 + 2x_2$ subject to three constraints. The constraints are written in a form where slack variables $w_1, w_2, w_3$ are explicitly defined. The variables $x_1, x_2, x_3$ are non-negative, as are the slack variables. The initial step involves setting $x_1$ and $x_2$ to zero, resulting in $\zeta = 0$. The notes then describe how to improve the solution by increasing $x_1$ while keeping $x_2 = 0$, based on the maximum coefficient rule. The constraints are analyzed to determine the maximum value of $x_1$ that maintains feasibility, which is found to be 5. Setting $x_1 = 5$ results in $w_3 = 0$, leading to a new vertex A with coordinates $(x_2 = 0, w_3 = 0)$. The notes then describe how to interchange $x_1$ and $w_3$, with $w_3$ leaving the basis and $x_1$ entering. The constraints are rewritten with $x_1$ expressed in terms of $w_3$ and $x_2$, and the objective function is updated accordingly.}%
\end{subfigure}%
\end{figure}

%
\newpage%
\section{Slide 10}%
\label{sec:Slide10}%
\subsection{Content}%

\circled{Ⅱ} \quad A \left( x_2 = 0, w_3 = 0 \right) \\
\underline{\text{Interchange } x_1 \text{ and } w_3 \text{ with }} \\
\underline{w_3 \text{ leaving basic and } x_1 \text{ entering basic}} \\
\text{max } \zeta = 15 - \frac{3}{2} w_3 + \frac{7}{2} x_2 \\
\text{subject to } \begin{cases}
x_1 = 5 - \frac{1}{2} w_3 + \frac{1}{2} x_2 \\
w_1 = 17 - \frac{1}{2} w_3 - \frac{5}{2} x_2 \\
w_2 = 3 + \frac{1}{2} w_3 - \frac{3}{2} x_2
\end{cases} \\
\underbrace{w_2}_{basic \ var.} \underbrace{x_2}_{non-basic \ var.}
%
\subsection{Description}%
Slide 10 of 44 in a Linear Programming course continues the simplex method iteration started in the previous slide.  The slide begins by stating that the current vertex is A, where $x_2 = 0$ and $w_3 = 0$.  It then describes the variable interchange step: $x_1$ enters the basis, and $w_3$ leaves the basis. This is clearly indicated by underlining. The updated objective function is given as $\text{max } \zeta = 15 - \frac{3}{2} w_3 + \frac{7}{2} x_2$. The constraints are rewritten to reflect this change, expressing $x_1$, $w_1$, and $w_2$ in terms of the new non-basic variables $w_3$ and $x_2$.  The resulting equations are: $x_1 = 5 - \frac{1}{2} w_3 + \frac{1}{2} x_2$, $w_1 = 17 - \frac{1}{2} w_3 - \frac{5}{2} x_2$, and $w_2 = 3 + \frac{1}{2} w_3 - \frac{3}{2} x_2$.  Finally, the slide labels the variables $w_2$ and $x_2$ as "basic var." and "non-basic var." respectively, indicating which variables are currently in and out of the basis. This sets the stage for the next iteration of the simplex method.%
\subsection{Figures}%
\begin{figure}[H]%
\centering%
\begin{subfigure}[b]{0.15\textwidth}%
\centering%

\includegraphics[width=\textwidth]{/Users/ashoksaravanan/Coding/ScribeLec/server/output/MA 421/lectures/2024-08-27-ExSimplex/figures/10.1.png}%
\caption{Handwritten notes describing the second iteration of the simplex method. The problem is to maximize $\zeta = 3x_1 + 2x_2$ subject to three constraints. The constraints are written in a form where slack variables $w_1, w_2, w_3$ are explicitly defined. The variables $x_1, x_2, x_3$ are non-negative, as are the slack variables. The initial step involves setting $x_1$ and $x_2$ to zero, resulting in $\zeta = 0$. The notes then describe how to improve the solution by increasing $x_1$ while keeping $x_2 = 0$, based on the maximum coefficient rule. The constraints are analyzed to determine the maximum value of $x_1$ that maintains feasibility, which is found to be 5. Setting $x_1 = 5$ results in $w_3 = 0$, leading to a new vertex A with coordinates $(x_2 = 0, w_3 = 0)$. The notes then describe how to interchange $x_1$ and $w_3$, with $w_3$ leaving the basis and $x_1$ entering. The constraints are rewritten with $x_1$ expressed in terms of $w_3$ and $x_2$, and the objective function is updated accordingly. The basic and non-basic variables are identified under the rewritten constraints.}%
\end{subfigure}%
\end{figure}

%
\newpage%
\section{Slide 11}%
\label{sec:Slide11}%
\subsection{Content}%

\circled{Ⅲ} \quad \text{Improve: } A \left( x_2 = 0, w_3 = 0 \right) \\
\text{increase } x_2, \text{ keeping } w_3 = 0 \\
\text{max } \zeta = 15 - \frac{3}{2} w_3 + \frac{7}{2} x_2 \\
\text{subject to } \begin{cases}
x_1 = 5 - \frac{1}{2} w_3 + \frac{1}{2} x_2 \\
w_1 = 17 - \frac{1}{2} w_3 - \frac{5}{2} x_2 \\
w_2 = 3 + \frac{1}{2} w_3 - \frac{3}{2} x_2
\end{cases} \\
\underbrace{\text{basic var.}}_{\text{x1, w1, w2}} \quad \underbrace{\text{non-basic var.}}_{\text{x2, w3}} \\
x_2 \text{ can be as large as possible} \\
w_1: 17 - \frac{5}{2} x_2 \ge 0 \implies x_2 \le \frac{34}{5} = 6.8 \\
w_2: 3 - \frac{3}{2} x_2 \ge 0 \implies x_2 \le 2 \\
\text{Set } x_2 = 2 \implies w_2 = 0 \text{ and } w_3 = 0 \implies B
%
\subsection{Description}%
Slide 11 of 44 in a Linear Programming course shows the third iteration of the simplex method.  Building upon slide 10, the current vertex is A ($x_2 = 0, w_3 = 0$). The goal is to improve the solution by increasing $x_2$ while keeping $w_3 = 0$. The objective function from slide 10, $\zeta = 15 - \frac{3}{2}w_3 + \frac{7}{2}x_2$, is used. The constraints from slide 10 are also used to determine the maximum feasible value of $x_2$.  The inequalities $17 - \frac{5}{2}x_2 \ge 0$ and $3 - \frac{3}{2}x_2 \ge 0$ are derived from $w_1 \ge 0$ and $w_2 \ge 0$ respectively. These yield $x_2 \le \frac{34}{5} = 6.8$ and $x_2 \le 2$. The smaller value, $x_2 = 2$, is selected to maintain feasibility.  Substituting $x_2 = 2$ and $w_3 = 0$ into the constraints shows that $w_2$ becomes 0, indicating that a new vertex B is reached. The slide clearly identifies basic and non-basic variables at each step, showing the transition from vertex A to vertex B. The slide visually connects the algebraic steps of the simplex method with the geometric interpretation of moving from one vertex to another in the feasible region.%
\subsection{Figures}%
\begin{figure}[H]%
\centering%
\begin{subfigure}[b]{0.15\textwidth}%
\centering%
\includegraphics[width=\textwidth]{/Users/ashoksaravanan/Coding/ScribeLec/server/output/MA 421/lectures/2024-08-27-ExSimplex/figures/11.1.png}%
\caption{Handwritten notes describing the third iteration of the simplex method. The problem is to maximize $\zeta = 3x_1 + 2x_2$ subject to three constraints. The constraints are written in a form where slack variables $w_1, w_2, w_3$ are explicitly defined. The variables $x_1, x_2, x_3$ are non-negative, as are the slack variables. The initial step involves setting $x_1$ and $x_2$ to zero, resulting in $\zeta = 0$. The notes then describe how to improve the solution by increasing $x_1$ while keeping $x_2 = 0$, based on the maximum coefficient rule. The constraints are analyzed to determine the maximum value of $x_1$ that maintains feasibility, which is found to be 5. Setting $x_1 = 5$ results in $w_3 = 0$, leading to a new vertex A with coordinates $(x_2 = 0, w_3 = 0)$. The notes then describe how to interchange $x_1$ and $w_3$, with $w_3$ leaving the basis and $x_1$ entering. The constraints are rewritten with $x_1$ expressed in terms of $w_3$ and $x_2$, and the objective function is updated accordingly. The basic and non-basic variables are identified under the rewritten constraints. The notes then describe how to improve the solution by increasing $x_2$ while keeping $w_3 = 0$. The constraints are analyzed to determine the maximum value of $x_2$ that maintains feasibility, which is found to be 2. Setting $x_2 = 2$ results in $w_2 = 0$, leading to a new vertex B.}%
\end{subfigure}%
\end{figure}

%
\newpage%
\section{Slide 12}%
\label{sec:Slide12}%
\subsection{Content}%

\circled{Ⅲ} \quad B \left( w_2 = 0, w_3 = 0 \right) \\
\underline{\text{Interchange } x_2 \text{ and } w_2 \text{ with }} \\
\underline{w_2 \text{ leaving basic and } x_2 \text{ entering basic}} \\
\text{max } \zeta = 15 - \frac{3}{2} w_3 + \frac{7}{2} x_2 \\
\text{subject to } \begin{cases}
x_1 = 5 - \frac{1}{2} w_3 + \frac{1}{2} x_2 \\
w_1 = 17 - \frac{1}{2} w_3 - \frac{5}{2} x_2 \\
w_2 = 3 + \frac{1}{2} w_3 - \frac{3}{2} x_2
\end{cases} \\
x_2 = 2 + \frac{1}{3} w_3 - \frac{2}{3} w_2 \\
x_1 = 5 - \frac{1}{2} w_3 + \frac{1}{2} \left( 2 + \frac{1}{3} w_3 - \frac{2}{3} w_2 \right) = 6 - \frac{2}{3} w_3 - \frac{1}{3} w_2 \\
w_1 = 17 - \frac{1}{2} w_3 - \frac{5}{2} \left( 2 + \frac{1}{3} w_3 - \frac{2}{3} w_2 \right) = 12 - \frac{4}{3} w_3 + \frac{5}{3} w_2
%
\subsection{Description}%
Slide 12 of 44 continues the simplex method, building on slide 11.  The current vertex is B ($w_2 = 0, w_3 = 0$). The slide indicates an interchange of variables: $x_2$ enters the basis, and $w_2$ leaves.  The objective function from slide 11 is shown, $\zeta = 15 - \frac{3}{2}w_3 + \frac{7}{2}x_2$.  The key step is solving the equation $w_2 = 3 + \frac{1}{2}w_3 - \frac{3}{2}x_2$ for $x_2$, resulting in $x_2 = 2 + \frac{1}{3}w_3 - \frac{2}{3}w_2$. This expression for $x_2$ is substituted into the equations for $x_1$ and $w_1$ from slide 11. The resulting equations are: $x_1 = 6 - \frac{2}{3}w_3 - \frac{1}{3}w_2$ and $w_1 = 12 - \frac{4}{3}w_3 + \frac{5}{3}w_2$.  The slide clearly shows the algebraic manipulation involved in pivoting from one basic feasible solution to another, moving from vertex B to the next iteration's vertex.  The underlining emphasizes the variable interchange.%
\subsection{Figures}%
\begin{figure}[H]%
\centering%
\begin{subfigure}[b]{0.15\textwidth}%
\centering%
\includegraphics[width=\textwidth]{/Users/ashoksaravanan/Coding/ScribeLec/server/output/MA 421/lectures/2024-08-27-ExSimplex/figures/12.1.png}%
\caption{Handwritten notes describing the third iteration of the simplex method. The problem is to maximize $\zeta = 3x_1 + 2x_2$ subject to three constraints. The constraints are written in a form where slack variables $w_1, w_2, w_3$ are explicitly defined. The variables $x_1, x_2, x_3$ are non-negative, as are the slack variables. The initial step involves setting $x_1$ and $x_2$ to zero, resulting in $\zeta = 0$. The notes then describe how to improve the solution by increasing $x_1$ while keeping $x_2 = 0$, based on the maximum coefficient rule. The constraints are analyzed to determine the maximum value of $x_1$ that maintains feasibility, which is found to be 5. Setting $x_1 = 5$ results in $w_3 = 0$, leading to a new vertex A with coordinates $(x_2 = 0, w_3 = 0)$. The notes then describe how to interchange $x_1$ and $w_3$, with $w_3$ leaving the basis and $x_1$ entering. The constraints are rewritten with $x_1$ expressed in terms of $w_3$ and $x_2$, and the objective function is updated accordingly. The basic and non-basic variables are identified under the rewritten constraints. The notes then describe how to improve the solution by increasing $x_2$ while keeping $w_3 = 0$. The constraints are analyzed to determine the maximum value of $x_2$ that maintains feasibility, which is found to be 2. Setting $x_2 = 2$ results in $w_2 = 0$, leading to a new vertex B. The notes then describe how to interchange $x_2$ and $w_2$, with $w_2$ leaving the basis and $x_2$ entering. The constraints are rewritten with $x_2$ expressed in terms of $w_3$ and $w_2$, and the expressions for $x_1$ and $w_1$ are updated accordingly.}%
\end{subfigure}%
\end{figure}

%
\newpage%
\section{Slide 13}%
\label{sec:Slide13}%
\subsection{Content}%

\circled{Ⅲ} \quad B \left( w_2 = 0, w_3 = 0 \right) \\
\underline{\text{Interchange } x_2 \text{ and } w_2 \text{ with }} \\
\underline{w_2 \text{ leaving basic and } x_2 \text{ entering basic}} \\
\text{max } \zeta = 15 - \frac{3}{2} w_3 + \frac{7}{2} x_2 \\
\text{subject to } \begin{cases}
x_1 = 5 - \frac{1}{2} w_3 + \frac{1}{2} x_2 \\
w_1 = 17 - \frac{1}{2} w_3 - \frac{5}{2} x_2 \\
w_2 = 3 + \frac{1}{2} w_3 - \frac{3}{2} x_2
\end{cases} \\
\downarrow \\
\zeta = 15 - \frac{3}{2} w_3 + \frac{7}{2} \left( 2 + \frac{1}{3} w_3 - \frac{2}{3} w_2 \right) \\
= 22 - \frac{1}{3} w_3 - \frac{7}{3} w_2 \\
x_2 = 2 + \frac{1}{3} w_3 - \frac{2}{3} w_2
%
\subsection{Description}%
Slide 13 of 44 shows the continuation of the simplex method from slide 12.  The current vertex is B, where $w_2 = 0$ and $w_3 = 0$.  The slide describes the variable interchange step: $x_2$ enters the basis, and $w_2$ leaves. This is clearly indicated by underlining. The objective function from slide 12, $\zeta = 15 - \frac{3}{2}w_3 + \frac{7}{2}x_2$, is used. The crucial step is substituting the expression for $x_2$ derived in slide 12, $x_2 = 2 + \frac{1}{3}w_3 - \frac{2}{3}w_2$, into the objective function. This substitution results in a new objective function: $\zeta = 22 - \frac{1}{3}w_3 - \frac{7}{3}w_2$.  The slide highlights this substitution process with an arrow. The updated objective function shows that the value of $\zeta$ at vertex B is 22. Since the coefficients of $w_3$ and $w_2$ are negative, this indicates that the current solution is optimal. The slide demonstrates the algebraic manipulation involved in updating the objective function after a variable interchange in the simplex method.%
\subsection{Figures}%
\begin{figure}[H]%
\centering%
\begin{subfigure}[b]{0.15\textwidth}%
\centering%

\includegraphics[width=\textwidth]{/Users/ashoksaravanan/Coding/ScribeLec/server/output/MA 421/lectures/2024-08-27-ExSimplex/figures/13.1.png}%
\caption{Handwritten notes describing the third iteration of the simplex method. The problem is to maximize $\zeta = 3x_1 + 2x_2$ subject to three constraints. The constraints are written in a form where slack variables $w_1, w_2, w_3$ are explicitly defined. The variables $x_1, x_2, x_3$ are non-negative, as are the slack variables. The initial step involves setting $x_1$ and $x_2$ to zero, resulting in $\zeta = 0$. The notes then describe how to improve the solution by increasing $x_1$ while keeping $x_2 = 0$, based on the maximum coefficient rule. The constraints are analyzed to determine the maximum value of $x_1$ that maintains feasibility, which is found to be 5. Setting $x_1 = 5$ results in $w_3 = 0$, leading to a new vertex A with coordinates $(x_2 = 0, w_3 = 0)$. The notes then describe how to interchange $x_1$ and $w_3$, with $w_3$ leaving the basis and $x_1$ entering. The constraints are rewritten with $x_1$ expressed in terms of $w_3$ and $x_2$, and the objective function is updated accordingly. The basic and non-basic variables are identified under the rewritten constraints. The notes then describe how to improve the solution by increasing $x_2$ while keeping $w_3 = 0$. The constraints are analyzed to determine the maximum value of $x_2$ that maintains feasibility, which is found to be 2. Setting $x_2 = 2$ results in $w_2 = 0$, leading to a new vertex B. The notes then describe how to interchange $x_2$ and $w_2$, with $w_2$ leaving the basis and $x_2$ entering. The constraints are rewritten with $x_2$ expressed in terms of $w_3$ and $w_2$, and the expressions for $x_1$ and $w_1$ are updated accordingly. The objective function is updated by substituting the expression for $x_2$ into the previous objective function.}%
\end{subfigure}%
\end{figure}

%
\newpage%
\section{Slide 14}%
\label{sec:Slide14}%
\subsection{Content}%

\circled{Ⅲ} \quad B \left( w_2 = 0, w_3 = 0 \right) \\
\text{max } \zeta = 22 - \frac{1}{3} w_3 - \frac{7}{3} w_2 \\
\text{subject to } \begin{cases}
x_1 = 6 - \frac{2}{3} w_3 - \frac{1}{3} w_2 \\
w_1 = 12 - \frac{4}{3} w_3 + \frac{5}{3} w_2 \\
x_2 = 2 + \frac{1}{3} w_3 - \frac{2}{3} w_2
\end{cases} \\
\underbrace{x_1, x_2, w_1}_{\text{basic}} \quad \underbrace{w_2, w_3}_{\text{non-basic}} \\
x_1, x_2, w_1, w_2, w_3 \ge 0 \\
B \text{ is optimal !!} \quad \text{max } \zeta = 22 \\
\zeta = 22 - \frac{1}{3} w_3 - \frac{7}{3} w_2
%
\subsection{Description}%
Slide 14 of 44 presents the final iteration of the simplex method.  Building upon the previous slides, the current vertex is B, with $w_2 = 0$ and $w_3 = 0$. The objective function is now $\zeta = 22 - \frac{1}{3}w_3 - \frac{7}{3}w_2$. The constraints are rewritten with $x_1$, $x_2$, and $w_1$ as basic variables and $w_2$ and $w_3$ as non-basic variables.  Crucially, the coefficients of the non-basic variables $w_2$ and $w_3$ in the objective function are negative (-7/3 and -1/3 respectively). This indicates that increasing either $w_2$ or $w_3$ would decrease the objective function. Therefore, the current solution is optimal, with a maximum value of $\zeta = 22$. The slide explicitly states "B is optimal!!" and "max $\zeta = 22$," confirming the optimality of the solution at vertex B. The slide summarizes the entire simplex method process, showing the final optimal solution and its value.%
\subsection{Figures}%
\begin{figure}[H]%
\centering%
\begin{subfigure}[b]{0.15\textwidth}%
\centering%
\includegraphics[width=\textwidth]{/Users/ashoksaravanan/Coding/ScribeLec/server/output/MA 421/lectures/2024-08-27-ExSimplex/figures/14.1.png}%
\caption{Handwritten notes describing the third iteration of the simplex method. The problem is to maximize $\zeta = 3x_1 + 2x_2$ subject to three constraints. The constraints are written in a form where slack variables $w_1, w_2, w_3$ are explicitly defined. The variables $x_1, x_2, x_3$ are non-negative, as are the slack variables. The initial step involves setting $x_1$ and $x_2$ to zero, resulting in $\zeta = 0$. The notes then describe how to improve the solution by increasing $x_1$ while keeping $x_2 = 0$, based on the maximum coefficient rule. The constraints are analyzed to determine the maximum value of $x_1$ that maintains feasibility, which is found to be 5. Setting $x_1 = 5$ results in $w_3 = 0$, leading to a new vertex A with coordinates $(x_2 = 0, w_3 = 0)$. The notes then describe how to interchange $x_1$ and $w_3$, with $w_3$ leaving the basis and $x_1$ entering. The constraints are rewritten with $x_1$ expressed in terms of $w_3$ and $x_2$, and the objective function is updated accordingly. The basic and non-basic variables are identified under the rewritten constraints. The notes then describe how to improve the solution by increasing $x_2$ while keeping $w_3 = 0$. The constraints are analyzed to determine the maximum value of $x_2$ that maintains feasibility, which is found to be 2. Setting $x_2 = 2$ results in $w_2 = 0$, leading to a new vertex B. The notes then describe how to interchange $x_2$ and $w_2$, with $w_2$ leaving the basis and $x_2$ entering. The constraints are rewritten with $x_2$ expressed in terms of $w_3$ and $w_2$, and the expressions for $x_1$ and $w_1$ are updated accordingly. The objective function is updated by substituting the expression for $x_2$ into the previous objective function. The final objective function is $\zeta = 22 - \frac{1}{3}w_3 - \frac{7}{3}w_2$. Since the coefficients of $w_3$ and $w_2$ are negative, the current solution is optimal, with $\zeta = 22$.}%
\end{subfigure}%
\end{figure}

%
\newpage%
\section{Slide 15}%
\label{sec:Slide15}%
\subsection{Content}%

\circled{Ⅲ} \quad B \left( w_2 = 0, w_3 = 0 \right) \\
\text{max } \zeta = 22 - \frac{1}{3} w_3 - \frac{7}{3} w_2 \\
\text{subject to } \begin{cases}
x_1 = 6 - \frac{2}{3} w_3 - \frac{1}{3} w_2 \\
w_1 = 12 - \frac{4}{3} w_3 + \frac{5}{3} w_2 \\
x_2 = 2 + \frac{1}{3} w_3 - \frac{2}{3} w_2
\end{cases} \\
x_1, x_2, w_1, w_2, w_3 \ge 0 \\
\downarrow \\
B \text{ is optimal!!} \quad \text{max } \zeta = 22 \\
\zeta = 22 - \frac{1}{3} w_3 - \frac{7}{3} w_2 \le 22, = 22 \text{ at } w_2 = 0, w_3 = 0
%
\subsection{Description}%
Slide 15 of 44 presents the final result of the simplex method.  The optimal solution is found at vertex B, where $w_2 = 0$ and $w_3 = 0$. The objective function is $\zeta = 22 - \frac{1}{3}w_3 - \frac{7}{3}w_2$.  Since the coefficients of $w_3$ and $w_2$ are negative, increasing either would decrease $\zeta$, confirming optimality. The constraints are shown with $x_1$, $x_2$, and $w_1$ as basic variables and $w_2$ and $w_3$ as non-basic variables. The optimal value of the objective function is $\zeta = 22$.  The notation "B is optimal!!" emphasizes the conclusion. A red arrow visually connects the constraints and the final optimal solution. The non-negativity constraints on all variables are explicitly stated. The context is the final step of the simplex algorithm in a linear programming problem, demonstrating the identification of the optimal solution and its value.%
\subsection{Figures}%
\begin{figure}[H]%
\centering%
\begin{subfigure}[b]{0.15\textwidth}%
\centering%
\includegraphics[width=\textwidth]{/Users/ashoksaravanan/Coding/ScribeLec/server/output/MA 421/lectures/2024-08-27-ExSimplex/figures/15.1.png}%
\caption{Handwritten notes describing the final iteration of the simplex method. The problem is to maximize $\zeta = 3x_1 + 2x_2$ subject to three constraints. The constraints are written in a form where slack variables $w_1, w_2, w_3$ are explicitly defined. The variables $x_1, x_2, x_3$ are non-negative, as are the slack variables. The notes show the final objective function $\zeta = 22 - \frac{1}{3}w_3 - \frac{7}{3}w_2$, the final constraints, and the conclusion that the optimal solution is at vertex B, where $w_2 = 0$ and $w_3 = 0$, with a maximum value of $\zeta = 22$. A red arrow points from the constraints to the conclusion.}%
\end{subfigure}%
\end{figure}

%
\newpage%
\section{Slide 16}%
\label{sec:Slide16}%
\subsection{Content}%

3x_1 + 2x_2 = 22 \\
x_1 = 0 \\
x_2 = 0 \\
D(x_1 = 0, w_1 \ne 0) \\
C(w_1 = 0, w_3 = 0) \\
-x_1 + 3x_2 = 12 \iff w_1 = 0 \\
B(w_2 = 0, w_3 = 0) \\
2x_1 - x_2 = 10 \iff w_3 = 0 \\
A(x_2 = 0, w_3 = 0) \\
x_1 + x_2 = 8 \iff w_2 = 0 \\
E(w_1 = 0, w_3 = 0) \text{ (not feasible)}
%
\subsection{Description}%
Slide 16 of 44 displays a graphical representation of the feasible region for a linear programming problem with two decision variables, $x_1$ and $x_2$, and three constraints.  The constraints are represented by lines: $-x_1 + 3x_2 = 12$ ($w_1 = 0$), $2x_1 - x_2 = 10$ ($w_3 = 0$), and $x_1 + x_2 = 8$ ($w_2 = 0$). The feasible region is shaded in light orange, representing the area satisfying all constraints.  The vertices of the feasible region are labeled: A ($x_2 = 0, w_3 = 0$), B ($w_2 = 0, w_3 = 0$), C ($w_1 = 0, w_3 = 0$), D ($x_1 = 0, w_1 \ne 0$), and E ($w_1 = 0, w_3 = 0$), which is explicitly marked as "not feasible".  The objective function, $3x_1 + 2x_2 = 22$, is shown as a line.  The optimal solution, found in the previous slides using the simplex method, is at vertex B, where the objective function line is tangent to the feasible region.  The Roman numerals I and II likely indicate the iterations of the simplex method leading to the optimal solution. The graph visually confirms the optimal solution obtained algebraically in the previous slides.%
\subsection{Figures}%
\begin{figure}[H]%
\centering%
\begin{subfigure}[b]{0.15\textwidth}%
\centering%
\includegraphics[width=\textwidth]{/Users/ashoksaravanan/Coding/ScribeLec/server/output/MA 421/lectures/2024-08-27-ExSimplex/figures/16.1.png}%
\caption{A graph showing the feasible region of a linear programming problem. The graph shows two variables, $x_1$ and $x_2$, and three constraints. The constraints are represented by lines, and the feasible region is the area where all constraints are satisfied. The vertices of the feasible region are labeled with their coordinates. The objective function is $3x_1 + 2x_2 = 22$, which is represented by a line. The optimal solution is at vertex B, where $w_2 = 0$ and $w_3 = 0$. The point E is labeled as not feasible.}%
\end{subfigure}%
\end{figure}

%
\newpage%
\section{Slide 17}%
\label{sec:Slide17}%
\subsection{Content}%

\text{[V] p. 11} \\
\text{max } \zeta = 5x_1 + 4x_2 + 3x_3 \\
\text{subject to } \begin{cases}
2x_1 + 3x_2 + x_3 \le 5 \\
4x_1 + x_2 + 2x_3 \le 11 \\
3x_1 + 4x_2 + 2x_3 \le 8 \\
x_1, x_2, x_3 \ge 0
\end{cases} \\
\overrightarrow{\text{red arrow}} \\
\begin{cases}
w_1 = 5 - 2x_1 - 3x_2 - x_3 \\
w_2 = 11 - 4x_1 - x_2 - 2x_3 \\
w_3 = 8 - 3x_1 - 4x_2 - 2x_3 \\
x_1, x_2, x_3, w_1, w_2, w_3 \ge 0
\end{cases}
%
\subsection{Description}%
Slide 17 of 44 introduces a linear programming problem. The objective is to maximize $\zeta = 5x_1 + 4x_2 + 3x_3$ subject to three inequality constraints: $2x_1 + 3x_2 + x_3 \le 5$, $4x_1 + x_2 + 2x_3 \le 11$, and $3x_1 + 4x_2 + 2x_3 \le 8$.  All variables ($x_1, x_2, x_3$) are non-negative.  The slide then shows the introduction of slack variables ($w_1, w_2, w_3$) to convert the inequalities into equalities.  The resulting system of equations is: $w_1 = 5 - 2x_1 - 3x_2 - x_3$, $w_2 = 11 - 4x_1 - x_2 - 2x_3$, and $w_3 = 8 - 3x_1 - 4x_2 - 2x_3$.  All variables, including the slack variables, are non-negative. A red arrow clearly indicates the transformation from the original problem to the equivalent problem with slack variables.  The notation "[V] p. 11" likely refers to a page number in a textbook or lecture notes.%
\subsection{Figures}%
\begin{figure}[H]%
\centering%
\begin{subfigure}[b]{0.15\textwidth}%
\centering%
\includegraphics[width=\textwidth]{/Users/ashoksaravanan/Coding/ScribeLec/server/output/MA 421/lectures/2024-08-27-ExSimplex/figures/17.1.png}%
\caption{Handwritten notes describing a linear programming problem. The problem is to maximize $\zeta = 5x_1 + 4x_2 + 3x_3$ subject to three constraints. The constraints are written in a form where slack variables $w_1, w_2, w_3$ are explicitly defined. The variables $x_1, x_2, x_3$ are non-negative, as are the slack variables. A red arrow points from the original constraints to the rewritten constraints with slack variables.}%
\end{subfigure}%
\end{figure}

%
\newpage%
\section{Slide 18}%
\label{sec:Slide18}%
\subsection{Content}%

\text{[V] p. 11} \quad \text{Set } x_1 = x_2 = x_3 = 0 \\
\circled{Ⅰ} \quad \text{max } \zeta = 5x_1 + 4x_2 + 3x_3 \\
\text{subject to } \begin{cases}
2x_1 + 3x_2 + x_3 \le 5 \\
4x_1 + x_2 + 2x_3 \le 11 \\
3x_1 + 4x_2 + 2x_3 \le 8 \\
x_1, x_2, x_3 \ge 0
\end{cases} \\
\overrightarrow{\text{red arrow}} \\
\begin{cases}
w_1 = 5 - 2x_1 - 3x_2 - x_3 = 5 \ge 0 \\
w_2 = 11 - 4x_1 - x_2 - 2x_3 = 11 \ge 0 \\
w_3 = 8 - 3x_1 - 4x_2 - 2x_3 = 8 \ge 0 \\
x_1, x_2, x_3, w_1, w_2, w_3 \ge 0
\end{cases}
%
\subsection{Description}%
Slide 18 of 44 presents a linear programming problem. The objective is to maximize $\zeta = 5x_1 + 4x_2 + 3x_3$ subject to three inequality constraints: $2x_1 + 3x_2 + x_3 \le 5$, $4x_1 + x_2 + 2x_3 \le 11$, and $3x_1 + 4x_2 + 2x_3 \le 8$.  All variables ($x_1, x_2, x_3$) are non-negative. The slide then introduces slack variables ($w_1, w_2, w_3$) to convert the inequalities into equalities, resulting in: $w_1 = 5 - 2x_1 - 3x_2 - x_3$, $w_2 = 11 - 4x_1 - x_2 - 2x_3$, and $w_3 = 8 - 3x_1 - 4x_2 - 2x_3$.  The non-negativity constraints are explicitly stated for all variables, including slack variables.  A red arrow visually connects the original constraints to their equivalent forms with slack variables.  The initial point considered is $x_1 = x_2 = x_3 = 0$, which is explicitly stated at the top of the slide. The notation "[V] p. 11" likely refers to a page number in a textbook or lecture notes. The circled Roman numeral I likely indicates this is the first iteration or a starting point in the simplex method.%
\subsection{Figures}%
\begin{figure}[H]%
\centering%
\begin{subfigure}[b]{0.15\textwidth}%
\centering%
\includegraphics[width=\textwidth]{/Users/ashoksaravanan/Coding/ScribeLec/server/output/MA 421/lectures/2024-08-27-ExSimplex/figures/18.1.png}%
\caption{Handwritten notes describing a linear programming problem. The problem is to maximize $\zeta = 5x_1 + 4x_2 + 3x_3$ subject to three constraints. The constraints are written in a form where slack variables $w_1, w_2, w_3$ are explicitly defined. The variables $x_1, x_2, x_3$ are non-negative, as are the slack variables. A red arrow points from the original constraints to the rewritten constraints with slack variables. The initial setting of $x_1=x_2=x_3=0$ is also specified.}%
\end{subfigure}%
\end{figure}

%
\newpage%
\section{Slide 19}%
\label{sec:Slide19}%
\subsection{Content}%

[V] p. 11 \quad \text{Set } x_1 = x_2 = x_3 = 0 \\
\text{max } \zeta = \boxed{5x_1} + 4x_2 + 3x_3 \\
\text{subject to } \begin{cases}
w_1 = 5 - 2x_1 - 3x_2 - x_3 \\
w_2 = 11 - 4x_1 - x_2 - 2x_3 \\
w_3 = 8 - 3x_1 - 4x_2 - 2x_3 \\
x_1, x_2, x_3, w_1, w_2, w_3 \ge 0
\end{cases} \\
\overrightarrow{\text{red arrow}} \\
\begin{cases}
w_3 \ge 0 \implies x_1 \le \frac{8}{3} \\
w_2 \ge 0 \implies x_1 \le \frac{11}{4} \\
w_1 \ge 0 \implies x_1 \le \frac{5}{2} \leftarrow \text{most strict}
\end{cases}
\text{increase } x_1, \text{ keep } x_2 = x_3 = 0
%
\subsection{Description}%
Slide 19 of 44 continues the solution of the linear programming problem introduced in the previous slides.  The problem is to maximize $\zeta = 5x_1 + 4x_2 + 3x_3$ subject to the constraints $2x_1 + 3x_2 + x_3 \le 5$, $4x_1 + x_2 + 2x_3 \le 11$, and $3x_1 + 4x_2 + 2x_3 \le 8$, with all variables non-negative.  The slide starts with the initial feasible solution $x_1 = x_2 = x_3 = 0$.  Slack variables $w_1, w_2, w_3$ have been introduced. The slide then focuses on increasing $x_1$ while keeping $x_2 = x_3 = 0$.  The constraints are rewritten to express the upper bounds on $x_1$ imposed by the non-negativity of the slack variables: $x_1 \le 8/3$, $x_1 \le 11/4$, and $x_1 \le 5/2$. The most restrictive constraint, $x_1 \le 5/2$, is identified as the one that will determine the next iteration's value of $x_1$.  The boxed $5x_1$ in the objective function highlights that $x_1$ is the variable to be increased in the next iteration of the simplex method. The red arrow indicates the flow of the analysis from the initial problem to the determination of the upper bound on $x_1$. The notation "[V] p. 11" likely refers to a page number in a textbook or lecture notes.%
\subsection{Figures}%
\begin{figure}[H]%
\centering%
\begin{subfigure}[b]{0.15\textwidth}%
\centering%
\includegraphics[width=\textwidth]{/Users/ashoksaravanan/Coding/ScribeLec/server/output/MA 421/lectures/2024-08-27-ExSimplex/figures/19.1.png}%
\caption{Handwritten notes describing a linear programming problem. The problem is to maximize $\zeta = 5x_1 + 4x_2 + 3x_3$ subject to three constraints. The constraints are written in a form where slack variables $w_1, w_2, w_3$ are explicitly defined. The variables $x_1, x_2, x_3$ are non-negative, as are the slack variables. A red arrow points from the original constraints to the rewritten constraints with slack variables. The initial setting of $x_1=x_2=x_3=0$ is also specified.  The constraints are then analyzed to determine the maximum value of $x_1$ while keeping $x_2=x_3=0$ and maintaining non-negativity of the slack variables.  A red arrow indicates the direction of the analysis. The most strict constraint is highlighted.}%
\end{subfigure}%
\end{figure}

%
\newpage%
\section{Slide 20}%
\label{sec:Slide20}%
\subsection{Content}%

[V] p. 11 \quad \text{Set } x_1 = x_2 = x_3 = 0 \\
\text{max } \zeta = \boxed{5x_1} + 4x_2 + 3x_3 \\
\text{subject to } \begin{cases}
w_1 = 5 - 2x_1 - 3x_2 - x_3 \\
w_2 = 11 - 4x_1 - x_2 - 2x_3 \\
w_3 = 8 - 3x_1 - 4x_2 - 2x_3 \\
x_1, x_2, x_3, w_1, w_2, w_3 \ge 0
\end{cases} \\
\overrightarrow{\text{red arrow}} \\
\text{increase } x_1, \text{ keep } x_2 = x_3 = 0 \\
\text{Set } x_1 = \frac{5}{2} \implies w_1 = 0, \text{ interchange } x_1, w_1 \\
\begin{cases}
x_1 = \frac{5}{2} - \frac{1}{2}w_1 - \frac{3}{2}x_2 - \frac{1}{2}x_3 \\
w_2 = 1 + 2w_1 + 5x_2 \\
w_3 = \frac{1}{2} + \frac{3}{2}w_1 + \frac{5}{2}x_2 - \frac{1}{2}x_3
\end{cases}
%
\subsection{Description}%
Slide 20 of 44 shows an iteration of the simplex method applied to the linear programming problem from previous slides. The objective is to maximize $\zeta = 5x_1 + 4x_2 + 3x_3$ subject to constraints involving slack variables $w_1, w_2, w_3$. Starting from $x_1 = x_2 = x_3 = 0$, the slide focuses on increasing $x_1$ while keeping $x_2 = x_3 = 0$.  The analysis from the previous slide leads to the conclusion that the maximum value for $x_1$ is $5/2$, which makes $w_1 = 0$.  This indicates that $w_1$ will leave the basis and $x_1$ will enter. The slide then shows the process of interchanging $x_1$ and $w_1$, rewriting the constraints to express $x_1$ as a function of $w_1, x_2, x_3$ and expressing $w_2$ and $w_3$ in terms of $w_1, x_2, x_3$. The boxed $5x_1$ in the objective function on the previous slide highlights that $x_1$ is the entering variable. The red arrow indicates the flow of the analysis. The notation "[V] p. 11" likely refers to a page number in a textbook or lecture notes. The context is clearly the simplex method for solving linear programming problems.%
\subsection{Figures}%
\begin{figure}[H]%
\centering%
\begin{subfigure}[b]{0.15\textwidth}%
\centering%

\includegraphics[width=\textwidth]{/Users/ashoksaravanan/Coding/ScribeLec/server/output/MA 421/lectures/2024-08-27-ExSimplex/figures/20.1.png}%
\caption{Handwritten notes describing a linear programming problem. The problem is to maximize $\zeta = 5x_1 + 4x_2 + 3x_3$ subject to three constraints. The constraints are written in a form where slack variables $w_1, w_2, w_3$ are explicitly defined. The variables $x_1, x_2, x_3$ are non-negative, as are the slack variables. A red arrow points from the original constraints to the rewritten constraints with slack variables. The initial setting of $x_1=x_2=x_3=0$ is also specified.  The constraints are then analyzed to determine the maximum value of $x_1$ while keeping $x_2=x_3=0$ and maintaining non-negativity of the slack variables.  The most strict constraint is used to set $x_1 = 5/2$, resulting in $w_1=0$.  The equations are then rewritten to express $x_1$ in terms of $w_1, x_2, x_3$ and the other slack variables in terms of $w_1, x_2, x_3$. A red arrow indicates the direction of the analysis. The most strict constraint is highlighted.  The process of interchanging $x_1$ and $w_1$ is explicitly noted.}%
\end{subfigure}%
\end{figure}

%
\newpage%
\section{Slide 21}%
\label{sec:Slide21}%
\subsection{Content}%

[V] p. 11 \quad \text{Set } x_1 = x_2 = x_3 = 0 \\
\text{max } \zeta = \boxed{5x_1} + 4x_2 + 3x_3 \quad \overrightarrow{\text{red arrow}} \quad \text{increase } x_1, \text{ keep } x_2 = x_3 = 0 \\
\text{subject to } \begin{cases}
w_1 = 5 - 2x_1 - 3x_2 - x_3 \\
w_2 = 11 - 4x_1 - x_2 - 2x_3 \\
w_3 = 8 - 3x_1 - 4x_2 - 2x_3 \\
x_1, x_2, x_3, w_1, w_2, w_3 \ge 0
\end{cases} \\
\text{Set } x_1 = \frac{5}{2} \implies w_1 = 0, \text{ interchange } x_1, w_1 \\
\zeta = 3x_1 + 4x_2 + 3x_3 \quad \overrightarrow{\text{red arrow}} \quad \zeta = \frac{25}{2} - \frac{5}{2}w_1 - \frac{7}{2}x_2 + \frac{1}{2}x_3
%
\subsection{Description}%
Slide 21 of 44 demonstrates a step in the simplex method for solving a linear programming problem.  The objective is to maximize $\zeta = 5x_1 + 4x_2 + 3x_3$ subject to three inequality constraints, which have been converted to equalities using slack variables $w_1, w_2, w_3$.  The initial feasible solution is $x_1 = x_2 = x_3 = 0$.  The slide shows the process of choosing the entering variable ($x_1$) and determining its value ($x_1 = 5/2$) based on the most restrictive constraint ($w_1 = 0$).  The variable $x_1$ enters the basis, and $w_1$ leaves.  The key step is rewriting the objective function $\zeta$ to express it in terms of the new basis, which now includes $x_1$ instead of $w_1$. The rewritten objective function is $\zeta = \frac{25}{2} - \frac{5}{2}w_1 - \frac{7}{2}x_2 + \frac{1}{2}x_3$. The red arrows highlight the transformation of the objective function. The notation "[V] p. 11" likely refers to a page number in a textbook or lecture notes. The boxed $5x_1$ in the objective function on the previous slide highlights that $x_1$ is the entering variable.  The context is clearly the simplex method for solving linear programming problems.%
\subsection{Figures}%
\begin{figure}[H]%
\centering%
\begin{subfigure}[b]{0.15\textwidth}%
\centering%
\includegraphics[width=\textwidth]{/Users/ashoksaravanan/Coding/ScribeLec/server/output/MA 421/lectures/2024-08-27-ExSimplex/figures/21.1.png}%
\caption{Handwritten notes describing a linear programming problem. The problem is to maximize $\zeta = 5x_1 + 4x_2 + 3x_3$ subject to three constraints. The constraints are written in a form where slack variables $w_1, w_2, w_3$ are explicitly defined. The variables $x_1, x_2, x_3$ are non-negative, as are the slack variables. A red arrow points from the original objective function to a rewritten objective function where $x_1$ is replaced using the equation for $w_1$. The initial setting of $x_1=x_2=x_3=0$ is also specified. The constraints are then analyzed to determine the maximum value of $x_1$ while keeping $x_2=x_3=0$ and maintaining non-negativity of the slack variables. The most strict constraint is used to set $x_1 = 5/2$, resulting in $w_1=0$. The equations are then rewritten to express $x_1$ in terms of $w_1, x_2, x_3$ and the other slack variables in terms of $w_1, x_2, x_3$. A red arrow indicates the direction of the analysis. The most strict constraint is highlighted. The process of interchanging $x_1$ and $w_1$ is explicitly noted. The objective function is rewritten to reflect the change of basis.}%
\end{subfigure}%
\end{figure}

%
\newpage%
\section{Slide 22}%
\label{sec:Slide22}%
\subsection{Content}%

[V] p. 11 \quad \text{Set } w_1 = 0, x_2 = 0, x_3 = 0 \\
\text{max. } \zeta = \frac{25}{2} - \frac{5}{2}w_1 - \frac{7}{2}x_2 + \frac{1}{2}x_3 \\
\text{s.t.} \begin{cases}
x_1 = \frac{5}{2} - \frac{1}{2}w_1 - \frac{3}{2}x_2 - \frac{1}{2}x_3 \\
w_2 = 1 + 2w_1 + 5x_2 \\
w_3 = \frac{1}{2} + \frac{3}{2}w_1 + \frac{5}{2}x_2 - \frac{1}{2}x_3
\end{cases} \\
\underbrace{\text{basic}}_{\text{basic}} \qquad \underbrace{\text{non-basic}}_{\text{non-basic}} \\
x_1, x_2, x_3, w_1, w_2, w_3 \ge 0
%
\subsection{Description}%
Slide 22 of 44 presents the result of the simplex method iteration from the previous slide. The objective function $\zeta$ has been rewritten in terms of the new basis, which now includes $x_1$ and excludes $w_1$. The constraints are also rewritten to express the basic variables ($x_1$, $w_2$, $w_3$) in terms of the non-basic variables ($w_1$, $x_2$, $x_3$).  The non-negativity constraints for all variables are explicitly stated. The terms "basic" and "non-basic" are clearly labeled to distinguish between the variables in the basis and those outside the basis. The initial values $w_1 = x_2 = x_3 = 0$ are specified, implying that the next step will involve determining which non-basic variable to bring into the basis. The context is clearly the simplex method for solving linear programming problems. The notation "[V] p. 11" likely refers to a page number in a textbook or lecture notes.%
\subsection{Figures}%
\begin{figure}[H]%
\centering%
\begin{subfigure}[b]{0.15\textwidth}%
\centering%

\includegraphics[width=\textwidth]{/Users/ashoksaravanan/Coding/ScribeLec/server/output/MA 421/lectures/2024-08-27-ExSimplex/figures/22.1.png}%
\caption{Handwritten notes showing the result of one iteration of the simplex method. The objective function is rewritten in terms of the new basis, which includes $x_1$ instead of $w_1$. The constraints are also rewritten to express the basic variables ($x_1, w_2, w_3$) in terms of the non-basic variables ($w_1, x_2, x_3$). The non-negativity constraints are explicitly stated. The words "basic" and "non-basic" are underlined and labeled to indicate which variables belong to each category. The initial values of $w_1, x_2, x_3$ are set to 0.}%
\end{subfigure}%
\end{figure}

%
\newpage%
\section{Slide 23}%
\label{sec:Slide23}%
\subsection{Content}%

[V] p. 11 \quad \text{Set } w_1 = 0, x_2 = 0, x_3 = 0 \\
\text{max. } \zeta = \frac{25}{2} - \frac{5}{2}w_1 - \frac{7}{2}x_2 + \frac{1}{2}x_3 \\
\text{s.t.} \begin{cases}
x_1 = \frac{5}{2} - \frac{1}{2}w_1 - \frac{3}{2}x_2 - \frac{1}{2}x_3 \\
w_2 = 1 + 2w_1 + 5x_2 \\
w_3 = \frac{1}{2} + \frac{3}{2}w_1 + \frac{5}{2}x_2 - \frac{1}{2}x_3
\end{cases} \\
\underbrace{\text{basic}}_{\text{basic}} \qquad \underbrace{\text{non-basic}}_{\text{non-basic}} \\
x_1, x_2, x_3, w_1, w_2, w_3 \ge 0 \\
\text{increase } x_3, \text{ keep } w_1, x_2 = 0 \\
w_3 \ge 0 \implies x_3 \le 1 \quad \leftarrow \text{most strict} \\
w_2 \ge 0 \implies x_3 \text{ can be any number} \\
x_1 \ge 0 \implies x_3 \le 5
%
\subsection{Description}%
Slide 23 of 44 continues the simplex method iteration.  The objective function and constraints from slide 22 are shown, with the basic variables ($x_1, w_2, w_3$) expressed in terms of the non-basic variables ($w_1, x_2, x_3$). The current solution has $w_1 = x_2 = x_3 = 0$, resulting in $x_1 = 5/2$, $w_2 = 1$, $w_3 = 1/2$, and $\zeta = 25/2$. The slide then focuses on increasing $x_3$ while keeping $w_1 = x_2 = 0$.  The constraints are checked to determine the maximum allowable increase in $x_3$ while maintaining non-negativity of the basic variables.  The analysis shows that $w_3 \ge 0$ implies $x_3 \le 1$, $w_2 \ge 0$ places no restriction on $x_3$, and $x_1 \ge 0$ implies $x_3 \le 5$. The most restrictive constraint, $x_3 \le 1$, is highlighted, indicating that $x_3$ will enter the basis and $w_3$ will leave in the next iteration. The context is clearly the simplex method for solving linear programming problems.%
\subsection{Figures}%
\begin{figure}[H]%
\centering%
\begin{subfigure}[b]{0.15\textwidth}%
\centering%
\includegraphics[width=\textwidth]{/Users/ashoksaravanan/Coding/ScribeLec/server/output/MA 421/lectures/2024-08-27-ExSimplex/figures/23.1.png}%
\caption{Handwritten notes showing the result of one iteration of the simplex method. The objective function is rewritten in terms of the new basis, which includes $x_1$ instead of $w_1$. The constraints are also rewritten to express the basic variables ($x_1, w_2, w_3$) in terms of the non-basic variables ($w_1, x_2, x_3$). The non-negativity constraints are explicitly stated. The words "basic" and "non-basic" are underlined and labeled to indicate which variables belong to each category. The initial values of $w_1, x_2, x_3$ are set to 0.  The next step is to increase $x_3$ while keeping $w_1$ and $x_2$ at 0. The constraints are analyzed to determine the maximum value of $x_3$ that maintains non-negativity of the basic variables. The most restrictive constraint is highlighted.}%
\end{subfigure}%
\end{figure}

%
\newpage%
\section{Slide 24}%
\label{sec:Slide24}%
\subsection{Content}%

[V] p. 11 \quad \text{Set } w_1 = 0, x_2 = 0, w_3 = 0 \\
\text{max. } \zeta = \frac{25}{2} - \frac{5}{2}w_1 - \frac{7}{2}x_2 + \frac{1}{2}x_3 \quad \begin{array}{c} \text{increase } x_3, \\ \text{keep } w_1, x_2 = 0 \end{array} \\
\text{s.t.} \begin{cases}
x_1 = \frac{5}{2} - \frac{1}{2}w_1 - \frac{3}{2}x_2 - \frac{1}{2}x_3 \\
w_2 = 1 + 2w_1 + 5x_2 \\
w_3 = \frac{1}{2} + \frac{3}{2}w_1 + \frac{5}{2}x_2 - \frac{1}{2}x_3
\end{cases} \\
\underbrace{\text{basic}}_{\text{basic}} \qquad \underbrace{\text{non-basic}}_{\text{non-basic}} \\
\text{Set } x_3 = 1 \implies w_3 = 0 \quad \text{Interchange } x_3, w_3 \\
x_3 = 1 + 3w_1 + x_2 - 2w_3 \\
x_1 = 2 - 2w_1 - 2x_2 + w_3 \\
w_2 = 1 + 2w_1 + 5x_2
%
\subsection{Description}%
Slide 24 of 44 shows another iteration of the simplex method.  The objective function and constraints from slide 23 are shown.  The previous slide determined that increasing $x_3$ while keeping $w_1 = x_2 = 0$ would lead to $w_3 = 0$ when $x_3 = 1$. This slide shows the result of pivoting, where $x_3$ enters the basis and $w_3$ leaves. The equations are rewritten to express the new basic variables ($x_1, x_3, w_2$) in terms of the non-basic variables ($w_1, x_2, w_3$).  The constraints are checked to determine the maximum allowable increase in $x_3$ while maintaining non-negativity of the basic variables. The most restrictive constraint is used to set $x_3 = 1$, resulting in $w_3 = 0$. The equations are then rewritten to express $x_3$ in terms of $w_1, x_2, w_3$ and the other basic variables in terms of $w_1, x_2, w_3$. The process of interchanging $x_3$ and $w_3$ is explicitly noted. The words "basic" and "non-basic" are underlined and labeled to indicate which variables belong to each category.  The red arrows highlight the flow of the calculations. The context is clearly the simplex method for solving linear programming problems.%
\subsection{Figures}%
\begin{figure}[H]%
\centering%
\begin{subfigure}[b]{0.15\textwidth}%
\centering%
\includegraphics[width=\textwidth]{/Users/ashoksaravanan/Coding/ScribeLec/server/output/MA 421/lectures/2024-08-27-ExSimplex/figures/24.1.png}%
\caption{Handwritten notes showing the continuation of the simplex method. The objective function and constraints from the previous slide are shown. The next step is to increase $x_3$ while keeping $w_1 = x_2 = 0$. The constraints are analyzed to determine the maximum allowable increase in $x_3$ while maintaining non-negativity of the basic variables. The most restrictive constraint is used to set $x_3 = 1$, resulting in $w_3 = 0$. The equations are then rewritten to express $x_3$ in terms of $w_1, x_2, w_3$ and the other basic variables in terms of $w_1, x_2, w_3$. The process of interchanging $x_3$ and $w_3$ is explicitly noted. The words "basic" and "non-basic" are underlined and labeled to indicate which variables belong to each category. A red arrow indicates the direction of the analysis. The most strict constraint is highlighted. The process of interchanging $x_3$ and $w_3$ is explicitly noted.}%
\end{subfigure}%
\end{figure}

%
\newpage%
\section{Slide 25}%
\label{sec:Slide25}%
\subsection{Content}%

[V] p. 11 \quad \text{Set } w_1 = 0, x_2 = 0, w_3 = 0 \\
\text{max. } \zeta = \frac{25}{2} - \frac{5}{2}w_1 - \frac{7}{2}x_2 + \frac{1}{2}x_3 \quad \begin{array}{c} \text{increase } x_3, \\ \text{keep } w_1, x_2 = 0 \end{array} \\
\text{s.t.} \begin{cases}
x_1 = \frac{5}{2} - \frac{1}{2}w_1 - \frac{3}{2}x_2 - \frac{1}{2}x_3 \\
w_2 = 1 + 2w_1 + 5x_2 \\
w_3 = \frac{1}{2} + \frac{3}{2}w_1 + \frac{5}{2}x_2 - \frac{1}{2}x_3
\end{cases} \\
\underbrace{\text{basic}}_{\text{basic}} \qquad \underbrace{\text{non-basic}}_{\text{non-basic}} \\
\text{Set } x_3 = 1 \implies w_3 = 0 \quad \text{Interchange } x_3, w_3 \\
\zeta = \frac{25}{2} - \frac{5}{2}w_1 - \frac{7}{2}x_2 + \frac{1}{2}x_3 \quad x_3 = 1 + 3w_1 + 5x_2 - 2w_3 \\
= 13 - w_1 - 3x_2 - w_3
%
\subsection{Description}%
Slide 25 of 44 continues the simplex algorithm.  It picks up where slide 24 left off, after determining that $x_3$ should enter the basis and $w_3$ should leave. The slide shows the substitution of $x_3 = 1$ into the objective function $\zeta$ and the simplification of the expression. The new objective function, expressed in terms of the new non-basic variables ($w_1, x_2, w_3$), is calculated. The constraints are rewritten to reflect the change in basis. The words "basic" and "non-basic" are used to label the variables. The red arrows and annotations guide the reader through the steps of the pivot operation. The context is the simplex method for solving linear programming problems. The notation "[V] p. 11" likely refers to a page number in a textbook or lecture notes.%
\subsection{Figures}%
\begin{figure}[H]%
\centering%
\begin{subfigure}[b]{0.15\textwidth}%
\centering%
\includegraphics[width=\textwidth]{/Users/ashoksaravanan/Coding/ScribeLec/server/output/MA 421/lectures/2024-08-27-ExSimplex/figures/25.1.png}%
\caption{Handwritten notes showing the continuation of the simplex method. The objective function and constraints from the previous slide are shown. The next step is to increase $x_3$ while keeping $w_1 = x_2 = 0$. The constraints are analyzed to determine the maximum allowable increase in $x_3$ while maintaining non-negativity of the basic variables. The most restrictive constraint is used to set $x_3 = 1$, resulting in $w_3 = 0$. The equations are then rewritten to express $x_3$ in terms of $w_1, x_2, w_3$ and the other basic variables in terms of $w_1, x_2, w_3$. The process of interchanging $x_3$ and $w_3$ is explicitly noted. The words "basic" and "non-basic" are underlined and labeled to indicate which variables belong to each category. A red arrow indicates the direction of the analysis. The most strict constraint is highlighted. The process of interchanging $x_3$ and $w_3$ is explicitly noted. The objective function is updated after the pivot operation.}%
\end{subfigure}%
\end{figure}

%
\newpage%
\section{Slide 26}%
\label{sec:Slide26}%
\subsection{Content}%

[V] p. 11 \quad \text{Set } w_1 = 0, x_2 = 0, w_3 = 0 \\
\text{max. } \zeta = \frac{25}{2} - \frac{5}{2}w_1 - \frac{7}{2}x_2 + \frac{1}{2}x_3 \quad \begin{array}{c} \text{increase } x_3, \\ \text{keep } w_1, x_2 = 0 \end{array} \\
\text{s.t.} \begin{cases}
x_1 = \frac{5}{2} - \frac{1}{2}w_1 - \frac{3}{2}x_2 - \frac{1}{2}x_3 \\
w_2 = 1 + 2w_1 + 5x_2 \\
w_3 = \frac{1}{2} + \frac{3}{2}w_1 + \frac{5}{2}x_2 - \frac{1}{2}x_3
\end{cases} \\
\underbrace{\text{basic}}_{\text{basic}} \qquad \underbrace{\text{non-basic}}_{\text{non-basic}} \\
\text{Set } x_3 = 1 \implies w_3 = 0 \quad \text{Interchange } x_3, w_3 \\
\zeta = 13 - w_1 - 3x_2 - w_3 \le 13 \\
\text{max } \zeta = 13, \text{ achieved at } w_1 = x_2 = w_3 = 0
%
\subsection{Description}%
Slide 26 of 44 presents the final step of a simplex method iteration.  Building upon the previous slides, it shows the updated objective function after the pivot operation where $x_3$ entered the basis and $w_3$ left. The objective function $\zeta$ is now expressed as $\zeta = 13 - w_1 - 3x_2 - w_3$.  Since all coefficients of the non-basic variables ($w_1, x_2, w_3$) are negative, the current solution ($w_1 = x_2 = w_3 = 0$) is optimal.  The maximum value of the objective function is 13. The words "basic" and "non-basic" are used to classify the variables.  The red arrows and annotations guide the reader through the steps. The notation "[V] p. 11" likely refers to a page number in a textbook or lecture notes. The context is clearly the simplex method for solving linear programming problems, specifically showing the optimal solution has been reached.%
\subsection{Figures}%
\begin{figure}[H]%
\centering%
\begin{subfigure}[b]{0.15\textwidth}%
\centering%

\includegraphics[width=\textwidth]{/Users/ashoksaravanan/Coding/ScribeLec/server/output/MA 421/lectures/2024-08-27-ExSimplex/figures/26.1.png}%
\caption{Handwritten notes showing the continuation of the simplex method. The objective function and constraints from the previous slide are shown. The next step is to increase $x_3$ while keeping $w_1 = x_2 = 0$. The constraints are analyzed to determine the maximum allowable increase in $x_3$ while maintaining non-negativity of the basic variables. The most restrictive constraint is used to set $x_3 = 1$, resulting in $w_3 = 0$. The equations are then rewritten to express $x_3$ in terms of $w_1, x_2, w_3$ and the other basic variables in terms of $w_1, x_2, w_3$. The process of interchanging $x_3$ and $w_3$ is explicitly noted. The words "basic" and "non-basic" are underlined and labeled to indicate which variables belong to each category. A red arrow indicates the direction of the analysis. The most strict constraint is highlighted. The process of interchanging $x_3$ and $w_3$ is explicitly noted. The objective function is updated after the pivot operation. The final result, showing the maximum value of the objective function and the values of the variables at the optimum, is presented.}%
\end{subfigure}%
\end{figure}

%
\newpage%
\section{Slide 27}%
\label{sec:Slide27}%
\subsection{Content}%

\text{Concept of a Dictionary of Variables} \\
\text{Initially } (x_1, x_2, \dots, x_n, w_1, w_2, \dots, w_m) \\
\underbrace{(x_1, x_2, \dots, x_n)}_{\text{n non-basic vars}} \quad \underbrace{(w_1, w_2, \dots, w_m)}_{\text{m basic variables}} \\
(x_1, x_2, \dots, x_n, x_{n+1}, x_{n+2}, \dots, x_{n+m}) \\
\begin{pmatrix}
x_{n+1} \\
x_{n+2} \\
\vdots \\
x_{n+m}
\end{pmatrix} = \begin{pmatrix}
b_1 \\
b_2 \\
\vdots \\
b_m
\end{pmatrix} - \begin{pmatrix}
a_{11} & a_{12} & \dots & a_{1n} \\
a_{21} & a_{22} & \dots & a_{2n} \\
\vdots & \vdots & \ddots & \vdots \\
a_{m1} & a_{m2} & \dots & a_{mn}
\end{pmatrix} \begin{pmatrix}
x_1 \\
x_2 \\
\vdots \\
x_n
\end{pmatrix} \\
\zeta = c_0 + c_1 x_1 + c_2 x_2 + \dots + c_n x_n
%
\subsection{Description}%
Slide 27 of 44 introduces the concept of a "dictionary of variables" within the context of linear programming.  The slide begins by defining an initial set of variables, separating them into $n$ non-basic variables ($x_1, x_2, \dots, x_n$) and $m$ basic variables ($w_1, w_2, \dots, w_m$).  A matrix equation is then presented, expressing the $m$ basic variables ($x_{n+1}, x_{n+2}, \dots, x_{n+m}$) as a linear combination of the $n$ non-basic variables ($x_1, x_2, \dots, x_n$) and a constant vector ($b_1, b_2, \dots, b_m$). The matrix $A = (a_{ij})$ represents the coefficients of the non-basic variables in the expressions for the basic variables. Finally, the objective function $\zeta$ is defined as a linear function of the non-basic variables. The red arrow visually connects the initial variable set to the matrix equation, illustrating how the dictionary represents the relationships between the variables. The underlining of "n non-basic vars" and "m basic variables" emphasizes the distinction between the two types of variables.  The context of linear programming is crucial for understanding the purpose of this dictionary, which is a key component of the simplex method for solving linear programming problems.%
\subsection{Figures}%
\begin{figure}[H]%
\centering%
\begin{subfigure}[b]{0.15\textwidth}%
\centering%
\includegraphics[width=\textwidth]{/Users/ashoksaravanan/Coding/ScribeLec/server/output/MA 421/lectures/2024-08-27-ExSimplex/figures/27.1.png}%
\caption{Handwritten notes describing the concept of a dictionary of variables in linear programming. The notes introduce the initial set of variables, classifying them as basic and non-basic. A matrix equation is presented, showing how the basic variables are expressed in terms of the non-basic variables and a constant vector. The objective function $\zeta$ is also defined. A red arrow indicates the relationship between the matrix equation and the initial set of variables. The words "n non-basic vars" and "m basic variables" are underlined.}%
\end{subfigure}%
\end{figure}

%
\newpage%
\section{Slide 28}%
\label{sec:Slide28}%
\subsection{Content}%

\text{Concept of a Dictionary of Variables} \\
\text{Initially } (x_1, x_2, \dots, x_n, w_1, w_2, \dots, w_m) \\
\underbrace{(x_1, x_2, \dots, x_n)}_{\text{n non-basic vars}} \quad \underbrace{(w_1, w_2, \dots, w_m)}_{\text{m basic variables}} \\
(x_1, x_2, \dots, x_n, x_{n+1}, x_{n+2}, \dots, x_{n+m}) \\
\underbrace{\curvearrowleft}_{\text{interchange one basic and one non-basic variable}}
%
\subsection{Description}%
Slide 28 of 44 introduces the concept of a "dictionary of variables" in linear programming, building upon the previous slide's simplex method illustration. It starts with an initial set of variables, divided into $n$ non-basic variables ($x_1, x_2, \dots, x_n$) and $m$ basic variables ($w_1, w_2, \dots, w_m$).  The slide then shows a second set of variables, where one basic variable has been exchanged with one non-basic variable. This exchange is visually represented by a red curved arrow, indicating the iterative nature of the simplex method. The terms "n non-basic vars" and "m basic variables" are underlined to emphasize the distinction. The slide sets the stage for a more formal representation of the simplex method using matrix notation, as seen in the previous slide. The context is crucial for understanding the iterative process of the simplex algorithm, where at each iteration, one basic and one non-basic variable are interchanged to improve the objective function.%
\subsection{Figures}%
\begin{figure}[H]%
\centering%
\begin{subfigure}[b]{0.15\textwidth}%
\centering%
\includegraphics[width=\textwidth]{/Users/ashoksaravanan/Coding/ScribeLec/server/output/MA 421/lectures/2024-08-27-ExSimplex/figures/28.1.png}%
\caption{Handwritten notes describing the concept of a dictionary of variables in linear programming. The notes introduce the initial set of variables, classifying them as basic and non-basic.  A red curved arrow indicates the process of interchanging one basic and one non-basic variable. The words "n non-basic vars" and "m basic variables" are underlined.}%
\end{subfigure}%
\end{figure}

%
\newpage%
\section{Slide 29}%
\label{sec:Slide29}%
\subsection{Content}%

\text{Concept of a Dictionary of Variables} \\
(\overline{x}_1, \overline{x}_2, \dots, \overline{x}_n, \overline{x}_{n+1}, \overline{x}_{n+2}, \dots, \overline{x}_{n+m}) \\
\underbrace{(\overline{x}_1, \overline{x}_2, \dots, \overline{x}_n)}_{n \text{ non-basic}} \quad \underbrace{(\overline{x}_{n+1}, \overline{x}_{n+2}, \dots, \overline{x}_{n+m})}_{m \text{ basic}} \\
\begin{pmatrix}
\overline{x}_{n+1} \\
\overline{x}_{n+2} \\
\vdots \\
\overline{x}_{n+m}
\end{pmatrix} = \begin{pmatrix}
\overline{b}_1 \\
\overline{b}_2 \\
\vdots \\
\overline{b}_m
\end{pmatrix} - \begin{pmatrix}
\overline{a}_{11} & \overline{a}_{12} & \dots & \overline{a}_{1n} \\
\overline{a}_{21} & \overline{a}_{22} & \dots & \overline{a}_{2n} \\
\vdots & \vdots & \ddots & \vdots \\
\overline{a}_{m1} & \overline{a}_{m2} & \dots & \overline{a}_{mn}
\end{pmatrix} \begin{pmatrix}
\overline{x}_1 \\
\overline{x}_2 \\
\vdots \\
\overline{x}_n
\end{pmatrix} \\
\overline{\zeta} = \overline{c}_0 + \overline{c}_1 \overline{x}_1 + \overline{c}_2 \overline{x}_2 + \dots + \overline{c}_n \overline{x}_n
%
\subsection{Description}%
Slide 29 of 44 formally defines the "dictionary of variables" used in the simplex method for linear programming.  It builds upon the previous slides by explicitly representing the relationships between basic and non-basic variables using matrix notation.  The slide begins by partitioning the variables into $n$ non-basic variables ($\overline{x}_1, \overline{x}_2, \dots, \overline{x}_n$) and $m$ basic variables ($\overline{x}_{n+1}, \overline{x}_{n+2}, \dots, \overline{x}_{n+m}$).  The core of the slide is a matrix equation expressing the basic variables as a linear function of the non-basic variables and a constant vector $\overline{b}$. The matrix $\overline{A} = (\overline{a}_{ij})$ contains the coefficients relating the non-basic and basic variables.  Finally, the objective function $\overline{\zeta}$ is defined as a linear function of the non-basic variables. The red dashed arrow visually connects the initial variable set to the matrix equation, emphasizing the representation of the system. The underlining of "n non-basic" and "m basic" highlights the classification of variables. The bars over the variables and coefficients indicate that these are the values at a particular iteration of the simplex method. The context is crucial because this slide provides the mathematical foundation for the iterative process of the simplex algorithm.%
\subsection{Figures}%
\begin{figure}[H]%
\centering%
\begin{subfigure}[b]{0.15\textwidth}%
\centering%
\includegraphics[width=\textwidth]{/Users/ashoksaravanan/Coding/ScribeLec/server/output/MA 421/lectures/2024-08-27-ExSimplex/figures/29.1.png}%
\caption{Handwritten notes describing the concept of a dictionary of variables in linear programming. The notes introduce the initial set of variables, classifying them as basic and non-basic. A matrix equation is presented, showing how the basic variables are expressed in terms of the non-basic variables and a constant vector. The objective function $\zeta$ is also defined. A red dashed arrow indicates the relationship between the matrix equation and the initial set of variables. The words "n non-basic" and "m basic" are underlined.}%
\end{subfigure}%
\end{figure}

%
\newpage%
\section{Slide 30}%
\label{sec:Slide30}%
\subsection{Content}%

\text{Concept of a Dictionary of Variables} \\
(1) \text{ In each dictionary, setting } n \text{ non-basic} \\
\underline{\text{variables to zero, corresponds to a}} \\
\underline{\text{vertex in a polyhedron.}} \\
(2) \text{ Total number of dictionaries} \\
= \binom{n+m}{n} = \frac{(n+m)!}{n! m!} \\
= \text{"\# of vertices"} \\
\text{(Note: not all vertices are feasible)}
%
\subsection{Description}%
Slide 30 of 44 discusses the relationship between dictionaries of variables in the simplex method and the vertices of the feasible region in a linear programming problem.  The slide makes two key points: (1) Each dictionary, obtained by setting the $n$ non-basic variables to zero, corresponds to a vertex of the feasible region (a polyhedron).  This is highlighted by underlining. (2) The total number of possible dictionaries is given by the binomial coefficient $\binom{n+m}{n}$, which is equivalent to $\frac{(n+m)!}{n!m!}$, where $n$ is the number of non-basic variables and $m$ is the number of basic variables. This is interpreted as the number of vertices of the polyhedron.  Importantly, the slide notes that not all vertices are feasible solutions to the linear program.  The context of linear programming is crucial here, as the slide connects the algebraic representation of the simplex method (dictionaries) to the geometric representation of the problem (vertices of the feasible region).%
\subsection{Figures}%
\begin{figure}[H]%
\centering%
\begin{subfigure}[b]{0.15\textwidth}%
\centering%

\includegraphics[width=\textwidth]{/Users/ashoksaravanan/Coding/ScribeLec/server/output/MA 421/lectures/2024-08-27-ExSimplex/figures/30.1.png}%
\caption{Handwritten notes describing the concept of a dictionary of variables in linear programming. The notes state that in each dictionary, setting the n non-basic variables to zero corresponds to a vertex in a polyhedron.  The total number of dictionaries is given by the binomial coefficient $\binom{n+m}{n}$, which is equivalent to $\frac{(n+m)!}{n!m!}$. This is interpreted as the number of vertices. A note is added that not all vertices are feasible.}%
\end{subfigure}%
\end{figure}

%
\newpage%
\section{Slide 31}%
\label{sec:Slide31}%
\subsection{Content}%

\text{What if the origin is not feasible?} \\
\text{[V] p.18} \\
\text{max } \zeta = -2x_1 - x_2 \\
\text{Subject to } -x_1 + x_2 \le -1 \\
\quad \quad \quad -x_1 - 2x_2 \le -2 \\
\quad \quad \quad x_2 \le 1 \\
\quad \quad \quad x_1, x_2 \ge 0
%
\subsection{Description}%
Slide 31 of 44 addresses a scenario in linear programming where the origin $(0,0)$ is not within the feasible region.  The slide presents a linear programming problem: maximize $\zeta = -2x_1 - x_2$ subject to the constraints $-x_1 + x_2 \le -1$, $-x_1 - 2x_2 \le -2$, $x_2 \le 1$, and $x_1, x_2 \ge 0$. A graph is included to visually represent the feasible region. The feasible region is a polygon in the first quadrant, bounded by the lines $-x_1 + x_2 = -1$, $-x_1 - 2x_2 = -2$, and $x_2 = 1$.  Crucially, the origin $(0,0)$ lies outside this feasible region, indicating that the constraints do not allow for a solution where both $x_1$ and $x_2$ are zero. The feasible region is shaded in orange, and the point (0,0) is marked with a red dot. The reference "[V] p.18" likely points to a specific page in a textbook or other reference material. The context of the slide is to introduce a situation where the standard simplex method, which typically starts at the origin, cannot be directly applied.  The slide sets the stage for discussing alternative methods to handle such cases.%
\subsection{Figures}%
\begin{figure}[H]%
\centering%
\begin{subfigure}[b]{0.15\textwidth}%
\centering%
\includegraphics[width=\textwidth]{/Users/ashoksaravanan/Coding/ScribeLec/server/output/MA 421/lectures/2024-08-27-ExSimplex/figures/31.1.png}%
\caption{A graph showing the feasible region of a linear programming problem. The feasible region is a polygon bounded by the lines $-x_1 + x_2 = -1$, $-x_1 - 2x_2 = -2$, and $x_2 = 1$. The origin $(0,0)$ is outside the feasible region. The feasible region is shaded in orange. The axes are labeled $x_1$ and $x_2$.}%
\end{subfigure}%
\end{figure}

%
\newpage%
\section{Slide 32}%
\label{sec:Slide32}%
\subsection{Content}%

\text{Auxiliary Problem} \\
\text{max } \xi = -x_0 \\
\text{subject to } -x_1 + x_2 - x_0 \le -1 \\
\quad \quad \quad -x_1 - 2x_2 - x_0 \le -2 \\
\quad \quad \quad x_2 - x_0 \le 1 \\
x_1, x_2, x_0 \ge 0 \\
(1) \text{ Must be feasible: choose } x_0 \text{ large enough} \\
(2) \text{ Original problem is feasible} \\
\iff \text{max } \xi = 0 \quad (\text{Note: max } \xi \le 0)
%
\subsection{Description}%
Slide 32 of 44 introduces the concept of an auxiliary problem in linear programming, used when the origin is not feasible in the original problem (as shown in the previous slide).  The auxiliary problem modifies the original problem by introducing a slack variable $x_0$. The objective function becomes maximizing $\xi = -x_0$, and the constraints are adjusted to include $x_0$. The constraints are modified to $-x_1 + x_2 - x_0 \le -1$, $-x_1 - 2x_2 - x_0 \le -2$, and $x_2 - x_0 \le 1$.  The key idea is that by choosing a sufficiently large $x_0$, the modified constraints become feasible. The slide then states that the original problem is feasible if and only if the maximum value of $\xi$ in the auxiliary problem is 0.  This is because if the original problem is feasible, there exists a solution where all original constraints are satisfied, and $x_0$ can be set to 0, resulting in $\xi = 0$. Conversely, if max $\xi = 0$, it implies that a feasible solution exists for the original problem. The note "max $\xi \le 0$" emphasizes that the maximum value of $\xi$ cannot be positive. The context is crucial because this slide presents a method to handle infeasible linear programming problems by transforming them into an equivalent feasible problem.%
\subsection{Figures}%
\begin{figure}[H]%
\centering%
\begin{subfigure}[b]{0.15\textwidth}%
\centering%
\includegraphics[width=\textwidth]{/Users/ashoksaravanan/Coding/ScribeLec/server/output/MA 421/lectures/2024-08-27-ExSimplex/figures/32.1.png}%
\caption{Handwritten notes describing the auxiliary problem in linear programming. The auxiliary problem introduces a new variable $x_0$ to ensure feasibility. The constraints are modified to include $x_0$, and the objective function is to maximize $\xi = -x_0$.  The notes explain that choosing a sufficiently large $x_0$ ensures the feasibility of the auxiliary problem.  The equivalence between the feasibility of the original problem and the condition max $\xi = 0$ is highlighted. A note indicates that the maximum value of $\xi$ is always less than or equal to 0.}%
\end{subfigure}%
\end{figure}

%
\newpage%
\section{Slide 33}%
\label{sec:Slide33}%
\subsection{Content}%

\text{Auxiliary Problem} \\
\text{max } \xi = -x_0 \\
\text{subject to } w_1 = -1 + x_1 - x_2 + x_0 \\
\quad \quad \quad w_2 = -2 + x_1 + 2x_2 + x_0 \\
\quad \quad \quad w_3 = 1 - x_2 + x_0 \\
\text{Setting } x_1, x_2, x_0 = 0 \\
w_1 = -1 \\
w_2 = -2 \leftarrow \text{most infeasible} \\
w_3 = 1 \quad \text{(interchange } x_0 \text{ and } w_2)
%
\subsection{Description}%
Slide 33 of 44 continues the discussion of the auxiliary problem introduced in slide 32.  This slide shows the application of the auxiliary problem to the specific linear program from slide 31. The auxiliary problem is formulated as: maximize $\xi = -x_0$ subject to $w_1 = -1 + x_1 - x_2 + x_0$, $w_2 = -2 + x_1 + 2x_2 + x_0$, and $w_3 = 1 - x_2 + x_0$, where $w_i$ are slack variables.  The slide then demonstrates what happens when we set the non-basic variables $x_1$, $x_2$, and $x_0$ to zero. This results in $w_1 = -1$, $w_2 = -2$, and $w_3 = 1$.  The slide highlights that $w_2 = -2$ represents the most infeasible constraint because it violates the constraint by the largest amount.  The note "(interchange $x_0$ and $w_2$)" indicates the next step in the simplex method:  to pivot on the $x_0$ and $w_2$ variables to move towards a feasible solution.  The context is crucial because this slide shows the practical application of the auxiliary problem and the first step in solving it using the simplex method. The slide demonstrates how to identify the most infeasible constraint and suggests the next step in the algorithm.%
\subsection{Figures}%
\begin{figure}[H]%
\centering%
\begin{subfigure}[b]{0.15\textwidth}%
\centering%
\includegraphics[width=\textwidth]{/Users/ashoksaravanan/Coding/ScribeLec/server/output/MA 421/lectures/2024-08-27-ExSimplex/figures/33.1.png}%
\caption{Handwritten notes describing the auxiliary problem in linear programming. The auxiliary problem introduces a new variable $x_0$ to ensure feasibility. The constraints are modified to include $x_0$, and the objective function is to maximize $\xi = -x_0$. The notes show the constraints $w_1 = -1 + x_1 - x_2 + x_0$, $w_2 = -2 + x_1 + 2x_2 + x_0$, and $w_3 = 1 - x_2 + x_0$. Setting $x_1, x_2, x_0 = 0$ results in $w_1 = -1$, $w_2 = -2$, and $w_3 = 1$. The most infeasible constraint is highlighted, and a note indicates that $x_0$ and $w_2$ should be interchanged.}%
\end{subfigure}%
\end{figure}

%
\newpage%
\section{Slide 34}%
\label{sec:Slide34}%
\subsection{Content}%

\text{Auxiliary Problem} \\
\text{max } \xi = -x_0 \\
\text{subject to } w_1 = -1 + x_1 - x_2 + x_0 \\
\quad \quad \quad w_2 = -2 + x_1 + 2x_2 + x_0 \\
\quad \quad \quad w_3 = 1 - x_2 + x_0 \\
x_0 = 2 - x_1 - 2x_2 + w_2 \\
w_1 = -1 + x_1 - x_2 + 2 - x_1 - 2x_2 + w_2 \\
= 1 - 3x_2 + w_2 \\
w_3 = 1 - x_2 + 2 - x_1 - 2x_2 + w_2 \\
= 3 - x_1 - 3x_2 + w_2
%
\subsection{Description}%
Slide 34 of 44 continues the solution of the auxiliary problem introduced in the previous slides.  The slide begins by restating the auxiliary problem: maximize $\xi = -x_0$ subject to $w_1 = -1 + x_1 - x_2 + x_0$, $w_2 = -2 + x_1 + 2x_2 + x_0$, and $w_3 = 1 - x_2 + x_0$.  The key step on this slide is solving the second constraint for $x_0$: $x_0 = 2 - x_1 - 2x_2 + w_2$. This expression for $x_0$ is then substituted into the first and third constraints to eliminate $x_0$.  The resulting simplified constraints are $w_1 = 1 - 3x_2 + w_2$ and $w_3 = 3 - x_1 - 3x_2 + w_2$. This prepares the problem for the next iteration of the simplex method, where the most infeasible constraint (identified in the previous slide) is addressed by bringing a new variable into the basis. The handwritten nature and the step-by-step calculations highlight the iterative process of solving linear programming problems using the simplex method. The context of the course (Linear Programming) is crucial for understanding the techniques used to solve the auxiliary problem and the overall goal of finding a feasible solution to the original problem.%
\subsection{Figures}%
\begin{figure}[H]%
\centering%
\begin{subfigure}[b]{0.15\textwidth}%
\centering%
\includegraphics[width=\textwidth]{/Users/ashoksaravanan/Coding/ScribeLec/server/output/MA 421/lectures/2024-08-27-ExSimplex/figures/34.1.png}%
\caption{Handwritten notes showing the steps to solve the auxiliary problem. The auxiliary problem is defined, and then a substitution is made for $x_0$ using the second constraint. This substitution is then used to simplify the first and third constraints.}%
\end{subfigure}%
\end{figure}

%
\newpage%
\section{Slide 35}%
\label{sec:Slide35}%
\subsection{Content}%

\text{Auxiliary Problem} \\
\text{max } \xi = -2 + x_1 + 2x_2 - w_2 \\
\text{subject to } x_0 = 2 - x_1 - 2x_2 + w_2 \\
\quad \quad \quad w_1 = 1 - 3x_2 + w_2 \\
\quad \quad \quad w_3 = 3 - x_1 - 3x_2 + w_2 \\
\text{Set } x_1, x_2, w_2 = 0 \\
\text{feasible!!}
%
\subsection{Description}%
Slide 35 of 44 shows the result of applying one iteration of the simplex method to the auxiliary problem from the previous slides. The auxiliary problem, maximize $\xi = -2 + x_1 + 2x_2 - w_2$, subject to $x_0 = 2 - x_1 - 2x_2 + w_2$, $w_1 = 1 - 3x_2 + w_2$, and $w_3 = 3 - x_1 - 3x_2 + w_2$, is presented.  The slide highlights the substitution of $x_0$ from the previous slide's calculations.  The values of $w_1$, $w_2$, and $w_3$ are calculated by setting $x_1 = x_2 = w_2 = 0$, resulting in $w_1 = 1$, $w_3 = 3$, and $x_0 = 2$.  Crucially, all $w_i \ge 0$, indicating that a feasible solution has been reached. The exclamation marks emphasize that a feasible solution to the auxiliary problem has been found.  The context is crucial because this slide demonstrates a successful step in the simplex method applied to the auxiliary problem, moving from an infeasible solution to a feasible one. The next step would be to continue the simplex method to optimize the objective function.%
\subsection{Figures}%
\begin{figure}[H]%
\centering%
\begin{subfigure}[b]{0.15\textwidth}%
\centering%

\includegraphics[width=\textwidth]{/Users/ashoksaravanan/Coding/ScribeLec/server/output/MA 421/lectures/2024-08-27-ExSimplex/figures/35.1.png}%
\caption{Handwritten notes showing the solution to the auxiliary problem after one iteration of the simplex method. The auxiliary problem is defined, and the constraints are rewritten after substituting for $x_0$. The values of $w_1$, $w_2$, and $w_3$ are highlighted when $x_1$, $x_2$, and $w_2$ are set to 0. The note "feasible!!" indicates that the solution is now feasible.}%
\end{subfigure}%
\end{figure}

%
\newpage%
\section{Slide 36}%
\label{sec:Slide36}%
\subsection{Content}%

\text{Auxiliary Problem} \\
\text{max } \xi = -2 + x_1 + 2x_2 - w_2 \\
\text{subject to } x_0 = 2 - x_1 - 2x_2 + w_2 \\
\quad \quad \quad w_1 = 1 - 3x_2 + w_2 \\
\quad \quad \quad w_3 = 3 - x_1 - 3x_2 + w_2 \\
\text{Set } x_1, x_2, w_2 = 0 \\
\text{feasible!!} \\
\text{increase } x_2 \\
x_0: x_2 \le 1 \\
w_1: x_2 \le \frac{1}{3} \\
w_3: x_2 \le 1 \leftarrow \text{interchange } x_2, w_1
%
\subsection{Description}%
Slide 36 of 44 continues the solution of the auxiliary problem.  The slide shows the auxiliary problem after the substitution of $x_0$ from the previous slide: maximize $\xi = -2 + x_1 + 2x_2 - w_2$ subject to $x_0 = 2 - x_1 - 2x_2 + w_2$, $w_1 = 1 - 3x_2 + w_2$, and $w_3 = 3 - x_1 - 3x_2 + w_2$.  Setting $x_1 = x_2 = w_2 = 0$ yields a feasible solution ($w_1 = 1$, $w_3 = 3$, $x_0 = 2$), as indicated by "feasible!!".  The next step, indicated by a red arrow and the words "increase $x_2$", involves increasing $x_2$ to improve the objective function.  The slide then calculates the upper bounds on $x_2$ imposed by each constraint: $x_2 \le 1$ from $x_0$, $x_2 \le \frac{1}{3}$ from $w_1$, and $x_2 \le 1$ from $w_3$.  The most restrictive constraint is $w_1$, and the note "interchange $x_2, w_1$" indicates that the next simplex iteration will involve pivoting on $x_2$ and $w_1$ to bring $x_2$ into the basis.  The context is crucial because this slide shows the iterative nature of the simplex method applied to the auxiliary problem, focusing on identifying the entering variable and the constraint that limits its increase.%
\subsection{Figures}%
\begin{figure}[H]%
\centering%
\begin{subfigure}[b]{0.15\textwidth}%
\centering%
\includegraphics[width=\textwidth]{/Users/ashoksaravanan/Coding/ScribeLec/server/output/MA 421/lectures/2024-08-27-ExSimplex/figures/36.1.png}%
\caption{Handwritten notes showing the solution to the auxiliary problem after one iteration of the simplex method. The auxiliary problem is defined, and the constraints are rewritten after substituting for $x_0$. The values of $w_1$, $w_2$, and $w_3$ are highlighted when $x_1$, $x_2$, and $w_2$ are set to 0. The note "feasible!!" indicates that the solution is now feasible.  A red arrow points to "increase $x_2$".  Below that, there are calculations showing the upper bounds on $x_2$ based on the constraints, and a note indicating that $x_2$ and $w_1$ should be interchanged.}%
\end{subfigure}%
\end{figure}

%
\newpage%
\section{Slide 37}%
\label{sec:Slide37}%
\subsection{Content}%

\text{Auxiliary Problem} \\
\text{max } \xi = -2 + x_1 + 2x_2 - w_2 \\
\text{subject to } x_0 = 2 - x_1 - 2x_2 + w_2 \\
\quad \quad \quad w_1 = 1 - 3x_2 + w_2 \\
\quad \quad \quad w_3 = 3 - x_1 - 3x_2 + w_2 \\
x_2 = \frac{1}{3} - \frac{1}{3}w_1 + \frac{1}{3}w_2 \\
w_3 = 2 - x_1 + w_1 \\
x_0 = \frac{4}{3} - \frac{1}{3}x_1 + \frac{2}{3}w_1 + \frac{1}{3}w_2
%
\subsection{Description}%
Slide 37 of 44 continues the solution of the auxiliary problem. The slide begins by restating the auxiliary problem: maximize $\xi = -2 + x_1 + 2x_2 - w_2$ subject to $x_0 = 2 - x_1 - 2x_2 + w_2$, $w_1 = 1 - 3x_2 + w_2$, and $w_3 = 3 - x_1 - 3x_2 + w_2$.  The key step on this slide is to perform further substitutions to express some variables in terms of others.  Specifically, the slide shows the calculation of $x_2$, $w_3$, and $x_0$ in terms of $x_1$, $w_1$, and $w_2$.  The numbers 1, 2, and 3 are highlighted in orange next to the constraints, possibly indicating the order of operations or the constraints being considered. Red arrows connect the constraints to the resulting expressions, showing the flow of the calculations. This step is part of the iterative process of the simplex method, aiming to find an optimal solution by systematically substituting variables and simplifying the constraints. The context of the course (Linear Programming) is crucial for understanding the techniques used to solve the auxiliary problem and the overall goal of finding a feasible solution to the original problem.%
\subsection{Figures}%
\begin{figure}[H]%
\centering%
\begin{subfigure}[b]{0.15\textwidth}%
\centering%
\includegraphics[width=\textwidth]{/Users/ashoksaravanan/Coding/ScribeLec/server/output/MA 421/lectures/2024-08-27-ExSimplex/figures/37.1.png}%
\caption{Handwritten notes showing the auxiliary problem and subsequent calculations. The auxiliary problem is defined, and then substitutions are made to express $x_2$, $w_3$, and $x_0$ in terms of other variables.  The numbers 1, 2, and 3 are highlighted in orange next to the constraints. Red arrows point from the constraints to the subsequent calculations.}%
\end{subfigure}%
\end{figure}

%
\newpage%
\section{Slide 38}%
\label{sec:Slide38}%
\subsection{Content}%

\text{Auxiliary Problem} \\
\text{max } \xi = -2 + x_1 + 2x_2 - w_2 \\
\text{subject to } x_0 = 2 - x_1 - 2x_2 + w_2 \\
\quad \quad \quad w_1 = 1 - 3x_2 + w_2 \\
\quad \quad \quad w_3 = 3 - x_1 - 3x_2 + w_2 \\
x_2 = \frac{1}{3} - \frac{1}{3}w_1 + \frac{1}{3}w_2 \\
\xi = -\frac{4}{3} + x_1 - \frac{2}{3}w_1 - \frac{1}{3}w_2
%
\subsection{Description}%
Slide 38 of 44 shows the continuation of the simplex method applied to the auxiliary problem.  The auxiliary problem, maximize $\xi = -2 + x_1 + 2x_2 - w_2$, subject to $x_0 = 2 - x_1 - 2x_2 + w_2$, $w_1 = 1 - 3x_2 + w_2$, and $w_3 = 3 - x_1 - 3x_2 + w_2$, is restated.  From the previous slide's analysis, it was determined that $x_2$ should enter the basis, and $w_1$ should leave.  This slide shows the result of this pivot operation.  The equation for $x_2$ is expressed in terms of $w_1$ and $w_2$: $x_2 = \frac{1}{3} - \frac{1}{3}w_1 + \frac{1}{3}w_2$.  This expression for $x_2$ is then substituted into the objective function $\xi$, resulting in a new expression for $\xi$ in terms of $x_1$, $w_1$, and $w_2$: $\xi = -\frac{4}{3} + x_1 - \frac{2}{3}w_1 - \frac{1}{3}w_2$. The numbers 1, 2, and 3 are highlighted in orange next to the constraints, possibly indicating the order of operations or the constraints being considered. Red arrows connect the constraints to the resulting expressions, showing the flow of the calculations. This step is part of the iterative process of the simplex method, aiming to find an optimal solution by systematically substituting variables and simplifying the constraints. The context of the course (Linear Programming) is crucial for understanding the techniques used to solve the auxiliary problem and the overall goal of finding a feasible solution to the original problem.%
\subsection{Figures}%
\begin{figure}[H]%
\centering%
\begin{subfigure}[b]{0.15\textwidth}%
\centering%

\includegraphics[width=\textwidth]{/Users/ashoksaravanan/Coding/ScribeLec/server/output/MA 421/lectures/2024-08-27-ExSimplex/figures/38.1.png}%
\caption{Handwritten notes showing the auxiliary problem and subsequent calculations. The auxiliary problem is defined, and then substitutions are made to express $x_2$ and $\xi$ in terms of other variables. The numbers 1, 2, and 3 are highlighted in orange next to the constraints. A red arrow points from the constraints to the calculation of $x_2$. Another red arrow points from the calculation of $x_2$ to the calculation of $\xi$.}%
\end{subfigure}%
\end{figure}

%
\newpage%
\section{Slide 39}%
\label{sec:Slide39}%
\subsection{Content}%

\text{Auxiliary Problem} \\
\text{max } \xi = -\frac{4}{3} + x_1 - \frac{2}{3}w_1 - \frac{1}{3}w_2 \\
\text{subject to } x_2 = \frac{1}{3} - \frac{1}{3}w_1 + \frac{1}{3}w_2 \\
\quad \quad \quad w_3 = 2 - x_1 + w_1 \\
\quad \quad \quad x_0 = \frac{4}{3} - x_1 + \frac{2}{3}w_1 + \frac{1}{3}w_2 \\
w_3: x_1 \le 2 \\
x_0: x_1 \le \frac{4}{3} \leftarrow \text{interchange } x_0, x_1 \\
\text{increase } x_1
%
\subsection{Description}%
Slide 39 of 44 shows another iteration of the simplex method applied to the auxiliary problem. The problem is to maximize $\xi = -\frac{4}{3} + x_1 - \frac{2}{3}w_1 - \frac{1}{3}w_2$ subject to $x_2 = \frac{1}{3} - \frac{1}{3}w_1 + \frac{1}{3}w_2$, $w_3 = 2 - x_1 + w_1$, and $x_0 = \frac{4}{3} - x_1 + \frac{2}{3}w_1 + \frac{1}{3}w_2$.  The objective function and constraints are rewritten after the previous pivot operation (Slide 38).  The next step involves increasing $x_1$ to further improve the objective function. The slide then calculates the upper bounds on $x_1$ imposed by each constraint: $x_1 \le 2$ from $w_3$, and $x_1 \le \frac{4}{3}$ from $x_0$. The most restrictive constraint is $x_0$, and the note "interchange $x_0, x_1$" indicates that the next simplex iteration will involve pivoting on $x_1$ and $x_0$ to bring $x_1$ into the basis.  The red arrow pointing to "increase $x_1$" clearly indicates the next step in the simplex algorithm. The context is crucial because this slide shows the iterative nature of the simplex method applied to the auxiliary problem, focusing on identifying the entering variable and the constraint that limits its increase.  The Roman numeral II next to the problem statement likely indicates this is the second iteration of the simplex method on the auxiliary problem.%
\subsection{Figures}%
\begin{figure}[H]%
\centering%
\begin{subfigure}[b]{0.15\textwidth}%
\centering%
\includegraphics[width=\textwidth]{/Users/ashoksaravanan/Coding/ScribeLec/server/output/MA 421/lectures/2024-08-27-ExSimplex/figures/39.1.png}%
\caption{Handwritten notes showing the auxiliary problem and subsequent calculations. The auxiliary problem is defined, and then substitutions are made to express $x_2$, $w_3$, and $x_0$ in terms of other variables. The numbers 1, 2, and 3 are highlighted in orange next to the constraints. A red arrow points from the constraints to the calculation of $x_2$. Another red arrow points from the calculation of $x_2$ to the calculation of $\xi$. A red arrow points to "increase $x_1$". Below that, there are calculations showing the upper bounds on $x_1$ based on the constraints, and a note indicating that $x_1$ and $x_0$ should be interchanged.}%
\end{subfigure}%
\end{figure}

%
\newpage%
\section{Slide 40}%
\label{sec:Slide40}%
\subsection{Content}%

\text{Auxiliary Problem} \\
\text{II} \\
\text{max } \xi = -\frac{4}{3} + x_1 - \frac{2}{3}w_1 - \frac{1}{3}w_2 \\
\text{subject to } x_2 = \frac{1}{3} - \frac{1}{3}w_1 + \frac{1}{3}w_2 \\
\quad \quad \quad w_3 = 2 - x_1 + w_1 \\
\quad \quad \quad x_0 = \frac{4}{3} - x_1 + \frac{2}{3}w_1 + \frac{1}{3}w_2 \\
x_1 = \frac{4}{3} + \frac{2}{3}w_1 + \frac{1}{3}w_2 - x_0 \\
w_3 = \frac{2}{3} + \frac{1}{3}w_1 - \frac{1}{3}w_2 + x_0 \\
x_2 = \frac{1}{3} - \frac{1}{3}w_1 + \frac{1}{3}w_2
%
\subsection{Description}%
Slide 40 of 44 continues the solution of the auxiliary problem from the previous slides. The auxiliary problem, to maximize $\xi = -\frac{4}{3} + x_1 - \frac{2}{3}w_1 - \frac{1}{3}w_2$, subject to $x_2 = \frac{1}{3} - \frac{1}{3}w_1 + \frac{1}{3}w_2$, $w_3 = 2 - x_1 + w_1$, and $x_0 = \frac{4}{3} - x_1 + \frac{2}{3}w_1 + \frac{1}{3}w_2$, is restated.  The key step on this slide is to solve for $x_1$ from the $x_0$ constraint, obtaining $x_1 = \frac{4}{3} + \frac{2}{3}w_1 + \frac{1}{3}w_2 - x_0$. This expression for $x_1$ is then substituted into the constraints for $w_3$ and $x_2$, resulting in new expressions for $w_3$ and $x_2$ in terms of $w_1$, $w_2$, and $x_0$. The number II is highlighted in orange, likely indicating this is the second iteration of the simplex method applied to the auxiliary problem. Red arrows show the flow of calculations, connecting the original constraints to the derived expressions.  The boxed equation for $x_1$ highlights its importance in the subsequent calculations. The context of the course (Linear Programming) is crucial for understanding the techniques used to solve the auxiliary problem and the overall goal of finding a feasible solution to the original problem. The slide demonstrates a step in the iterative process of the simplex method, where variables are systematically substituted and constraints are rewritten to find an optimal solution.%
\subsection{Figures}%
\begin{figure}[H]%
\centering%
\begin{subfigure}[b]{0.15\textwidth}%
\centering%
\includegraphics[width=\textwidth]{/Users/ashoksaravanan/Coding/ScribeLec/server/output/MA 421/lectures/2024-08-27-ExSimplex/figures/40.1.png}%
\caption{Handwritten notes showing the auxiliary problem and subsequent calculations. The auxiliary problem is defined, and then substitutions are made to express $x_1$, $w_3$, and $x_2$ in terms of other variables. The number II is highlighted in orange next to the problem statement. Red arrows point from the constraints to the subsequent calculations. A box is drawn around the equation for $x_1$.}%
\end{subfigure}%
\end{figure}

%
\newpage%
\section{Slide 41}%
\label{sec:Slide41}%
\subsection{Content}%

\text{Auxiliary Problem} \\
\text{II} \\
\text{max } \xi = -\frac{4}{3} + x_1 - \frac{2}{3}w_1 - \frac{1}{3}w_2 \\
\text{subject to } x_2 = \frac{1}{3} - \frac{1}{3}w_1 + \frac{1}{3}w_2 \\
\quad \quad \quad w_3 = 2 - x_1 + w_1 \\
\quad \quad \quad x_0 = \frac{4}{3} - x_1 + \frac{2}{3}w_1 + \frac{1}{3}w_2 \\
x_1 = \frac{4}{3} + \frac{2}{3}w_1 + \frac{1}{3}w_2 - x_0 \\
w_3 = \frac{2}{3} + \frac{1}{3}w_1 - \frac{1}{3}w_2 + x_0 \\
x_2 = \frac{1}{3} - \frac{1}{3}w_1 + \frac{1}{3}w_2
%
\subsection{Description}%
Slide 41 of 44 shows the continuation of the simplex method applied to the auxiliary problem.  The auxiliary problem, to maximize $\xi = -\frac{4}{3} + x_1 - \frac{2}{3}w_1 - \frac{1}{3}w_2$, subject to $x_2 = \frac{1}{3} - \frac{1}{3}w_1 + \frac{1}{3}w_2$, $w_3 = 2 - x_1 + w_1$, and $x_0 = \frac{4}{3} - x_1 + \frac{2}{3}w_1 + \frac{1}{3}w_2$, is restated.  From the previous slide, it was determined that $x_1$ should enter the basis and $x_0$ should leave. This slide shows the result of this pivot operation. The equation for $x_1$ is expressed in terms of $w_1$, $w_2$, and $x_0$: $x_1 = \frac{4}{3} + \frac{2}{3}w_1 + \frac{1}{3}w_2 - x_0$. This expression for $x_1$ is then substituted into the constraints for $w_3$ and $x_2$, resulting in new expressions for $w_3$ and $x_2$ in terms of $w_1$, $w_2$, and $x_0$. The number II is highlighted in orange, likely indicating this is the second iteration of the simplex method applied to the auxiliary problem. Red arrows show the flow of calculations, connecting the original constraints to the derived expressions. The boxed equation for $x_1$ highlights its importance in the subsequent calculations. The context of the course (Linear Programming) is crucial for understanding the techniques used to solve the auxiliary problem and the overall goal of finding a feasible solution to the original problem. The slide demonstrates a step in the iterative process of the simplex method, where variables are systematically substituted and constraints are rewritten to find an optimal solution.%
\subsection{Figures}%
\begin{figure}[H]%
\centering%
\begin{subfigure}[b]{0.15\textwidth}%
\centering%

\includegraphics[width=\textwidth]{/Users/ashoksaravanan/Coding/ScribeLec/server/output/MA 421/lectures/2024-08-27-ExSimplex/figures/41.1.png}%
\caption{Handwritten notes showing the auxiliary problem and subsequent calculations. The auxiliary problem is defined, and then substitutions are made to express $x_1$, $w_3$, and $x_2$ in terms of other variables. The number II is highlighted in orange next to the problem statement. Red arrows point from the constraints to the subsequent calculations. A box is drawn around the equation for $x_1$.}%
\end{subfigure}%
\end{figure}

%
\newpage%
\section{Slide 42}%
\label{sec:Slide42}%
\subsection{Content}%

\text{Auxiliary Problem} \\
\text{Ⅲ} \\
\text{max } \xi = -x_0 \\
\text{subject to } x_1 = \frac{4}{3} + \frac{2}{3}w_1 + \frac{1}{3}w_2 - x_0 \\
\quad \quad \quad w_3 = \frac{2}{3} + \frac{1}{3}w_1 - \frac{1}{3}w_2 + x_0 \\
\quad \quad \quad x_2 = \frac{1}{3} - \frac{1}{3}w_1 + \frac{1}{3}w_2 \\
\text{Setting } w_1 = w_2 = x_0 = 0 \\
\text{max } \xi = -x_0 = 0 \quad \text{Optimal!}
%
\subsection{Description}%
Slide 42 of 44 presents the final step in solving the auxiliary problem using the simplex method. The auxiliary problem, to maximize $\xi = -x_0$, subject to $x_1 = \frac{4}{3} + \frac{2}{3}w_1 + \frac{1}{3}w_2 - x_0$, $w_3 = \frac{2}{3} + \frac{1}{3}w_1 - \frac{1}{3}w_2 + x_0$, and $x_2 = \frac{1}{3} - \frac{1}{3}w_1 + \frac{1}{3}w_2$, is shown.  The solution is obtained by setting $w_1 = w_2 = x_0 = 0$. This results in $\xi = -x_0 = 0$, which is declared as the optimal solution. The Roman numeral III next to the problem statement likely indicates this is the third iteration of the simplex method on the auxiliary problem. The red arrow highlights the flow from the constraints to the final solution. The statement "Optimal!" confirms that a feasible solution to the original problem has been found.  The context of linear programming is crucial here because the auxiliary problem is a technique used to find an initial feasible solution for the original linear programming problem when the initial solution is not feasible.  The fact that the optimal value of the auxiliary problem is 0 indicates that a feasible solution to the original problem exists.%
\subsection{Figures}%
\begin{figure}[H]%
\centering%
\begin{subfigure}[b]{0.15\textwidth}%
\centering%
\includegraphics[width=\textwidth]{/Users/ashoksaravanan/Coding/ScribeLec/server/output/MA 421/lectures/2024-08-27-ExSimplex/figures/42.1.png}%
\caption{Handwritten notes showing the auxiliary problem and its solution. The auxiliary problem is defined, and then substitutions are made to express $x_1$, $w_3$, and $x_2$ in terms of other variables. The number III is highlighted in orange next to the problem statement. A red arrow points from the constraints to the calculation of $x_1$, $w_3$, and $x_2$. A red arrow points from the constraints to the final solution. Below that, there is a statement indicating that setting $w_1$, $w_2$, and $x_0$ to 0 yields an optimal solution of 0.}%
\end{subfigure}%
\end{figure}

%
\newpage%
\section{Slide 43}%
\label{sec:Slide43}%
\subsection{Content}%

\text{Back to Original Problem} \\
\text{max } \xi = -2x_1 - x_2 \\
\text{subject to } -x_1 + x_2 \le -1 \\
\quad \quad \quad -x_1 - 2x_2 \le -2 \\
\quad \quad \quad x_2 \le 1 \\
\quad \quad \quad x_1, x_2 \ge 0 \\
\text{We already have} \\
\begin{cases}
x_1 = \frac{4}{3} + \frac{2}{3}w_1 + \frac{1}{3}w_2 - x_0 \\
w_3 = \frac{2}{3} + \frac{1}{3}w_1 - \frac{1}{3}w_2 + x_0 \\
x_2 = \frac{1}{3} - \frac{1}{3}w_1 + \frac{1}{3}w_2
\end{cases}
%
\subsection{Description}%
Slide 43 of 44 shows the return to the original linear programming problem after solving the auxiliary problem. The original problem is restated: maximize $\xi = -2x_1 - x_2$ subject to the constraints $-x_1 + x_2 \le -1$, $-x_1 - 2x_2 \le -2$, $x_2 \le 1$, and $x_1, x_2 \ge 0$.  The solution to the auxiliary problem provides expressions for $x_1$, $x_2$, and $w_3$ in terms of $w_1$, $w_2$, and $x_0$.  Crucially, the solution to the auxiliary problem ($x_0 = 0$) is substituted into these expressions, effectively eliminating $x_0$ and providing a feasible solution to the original problem. The red strikethroughs on $x_0$ in the equations for $x_1$ and $w_3$ indicate this substitution. The red braces highlight the sets of equations for the original problem and the solution obtained from the auxiliary problem. The context of linear programming is essential here because the auxiliary problem is a tool to find an initial feasible solution for the original problem, which is then used to start the simplex method for the original problem.  The slide shows the transition from the auxiliary problem back to the original problem, using the solution of the auxiliary problem to obtain a starting point for solving the original problem.%
\subsection{Figures}%
\begin{figure}[H]%
\centering%
\begin{subfigure}[b]{0.15\textwidth}%
\centering%

\includegraphics[width=\textwidth]{/Users/ashoksaravanan/Coding/ScribeLec/server/output/MA 421/lectures/2024-08-27-ExSimplex/figures/43.1.png}%
\caption{Handwritten notes showing the original problem and the solution obtained from the auxiliary problem. The original problem is to maximize $\xi = -2x_1 - x_2$ subject to $-x_1 + x_2 \le -1$, $-x_1 - 2x_2 \le -2$, $x_2 \le 1$, and $x_1, x_2 \ge 0$. Below that, there is a set of equations for $x_1$, $w_3$, and $x_2$ in terms of $w_1$, $w_2$, and $x_0$, obtained from solving the auxiliary problem.  A red brace connects the constraints of the original problem. A red brace connects the equations obtained from the auxiliary problem.  A red strikethrough is present on $x_0$ in the equations for $x_1$ and $w_3$.}%
\end{subfigure}%
\end{figure}

%
\newpage%
\section{Slide 44}%
\label{sec:Slide44}%
\subsection{Content}%

\text{Back to Original Problem} \\
\text{max } \xi = -2x_1 - x_2 \\
\begin{cases}
x_1 = \frac{4}{3} + \frac{2}{3}w_1 + \frac{1}{3}w_2 \\
w_3 = \frac{2}{3} + \frac{1}{3}w_1 - \frac{1}{3}w_2 \\
x_2 = \frac{1}{3} - \frac{1}{3}w_1 + \frac{1}{3}w_2
\end{cases} \\
x_1, x_2, w_1, w_2, w_3 \ge 0 \\
\xi = -2x_1 - x_2 \\
= -3 - w_1 - w_2 \\
\text{Already at Optimal} \\
(w_1 = w_2 = 0)
%
\subsection{Description}%
Slide 44 of 44 presents the final solution to the original linear programming problem.  The original objective function, $\xi = -2x_1 - x_2$, is shown. The expressions for $x_1$, $x_2$, and $w_3$ obtained from the auxiliary problem (with $x_0 = 0$) are substituted into the objective function. This substitution leads to a simplified expression for $\xi$ in terms of $w_1$ and $w_2$: $\xi = -3 - w_1 - w_2$. Since $w_1$ and $w_2$ must be non-negative, the maximum value of $\xi$ is achieved when $w_1 = w_2 = 0$, resulting in $\xi = -3$.  The note "Already at Optimal" indicates that no further iterations of the simplex method are needed. The red arrows illustrate the flow of calculations, showing how the solution from the auxiliary problem is used to find the optimal solution to the original problem. The context of linear programming is crucial here because this slide demonstrates the final step in solving a linear programming problem using the two-phase simplex method. The auxiliary problem provided a feasible starting point, and the solution to the auxiliary problem is used to find the optimal solution to the original problem. The optimal solution is $\xi = -3$ when $w_1 = w_2 = 0$, which implies specific values for $x_1$ and $x_2$ based on the equations derived from the auxiliary problem.%
\subsection{Figures}%
\begin{figure}[H]%
\centering%
\begin{subfigure}[b]{0.15\textwidth}%
\centering%

\includegraphics[width=\textwidth]{/Users/ashoksaravanan/Coding/ScribeLec/server/output/MA 421/lectures/2024-08-27-ExSimplex/figures/44.1.png}%
\caption{Handwritten notes showing the original problem and the solution obtained from the auxiliary problem. The original problem is to maximize $\xi = -2x_1 - x_2$ subject to constraints (not explicitly written but implied from previous slides). Below that, there is a set of equations for $x_1$, $w_3$, and $x_2$ in terms of $w_1$ and $w_2$, obtained from solving the auxiliary problem.  The solution $x_1$, $x_2$ are substituted into the objective function $\xi = -2x_1 - x_2$, resulting in $\xi = -3 - w_1 - w_2$.  A red arrow points from the constraints to the calculation of $x_1$, $x_2$, and $w_3$. A red arrow points from the equations for $x_1$ and $x_2$ to the objective function.  A note indicates that the solution is already optimal when $w_1 = w_2 = 0$.}%
\end{subfigure}%
\end{figure}

%
\newpage%
\end{document}